% =====================================================================
% SUPPLEMENTARY MATERIAL — companion to atlas80030_main.tex
%
% Revision 2. The two twenty-column tables of the previous revision have been
% split into grouped, readable tables, and the paired SHAP-Delta table has been
% separated into its per-feature part and its systematic-level score and rank
% part. Every value is copied verbatim from the frozen artifacts by
% tables/make_supplementary_tables.py, which asserts that no numeric token is
% lost, changed or invented in the re-grouping; nothing is retyped here.
% =====================================================================
\documentclass[11pt,a4paper,onecolumn]{article}

\usepackage[margin=1.8cm]{geometry}
\usepackage{graphicx}
\usepackage{booktabs}
\usepackage{longtable}
\usepackage{pdflscape}
\usepackage{amsmath,amssymb}
\usepackage{url}
\newcommand{\pt}{p_{\mathrm{T}}}
\usepackage[colorlinks=true,linkcolor=blue,urlcolor=blue]{hyperref}

\renewcommand{\thetable}{S\arabic{table}}
\renewcommand{\thesection}{S\arabic{section}}

% Float placement: allow a full page of tables without forcing blank pages.
\renewcommand{\topfraction}{0.95}
\renewcommand{\bottomfraction}{0.95}
\renewcommand{\textfraction}{0.05}
\renewcommand{\floatpagefraction}{0.60}
\setcounter{topnumber}{3}
\setcounter{totalnumber}{4}

\title{Supplementary Material\\[4pt]
\large Systematic robustness of a boosted top-quark tagger on ATLAS open data:\\
a constituent-level feature reconstruction, paired SHAP-Delta attribution\\
drift, and conformal coverage}
\author{Adeel Rahmat \and M. Imran Jamil \and Ahmed Awan}
\date{}

\begin{document}
\maketitle

\section{Scope and provenance}

This document contains the complete result tables referred to in the main text.
Every value is taken from the frozen publication artifacts of a single analysis
chain, identified by

\begin{itemize}
  \item reconstruction fingerprint \texttt{c28a262540824e9db90f34c421f10c70}
        (810 cache files, 80{,}481{,}421 jet entries, seventeen samples);
  \item final XGBoost model SHA-256\\
        \texttt{63a23308efbf2ce90418c8e657730f57}\\
        \texttt{1db4798fc9ca2fb3692cb986de405f4c},
        developed from a five-million-jet nominal subset of which
        $4{,}499{,}612$ jets were used for fitting and $500{,}385$ were held out
        for validation and conformal calibration;
  \item metrics artifact
        \texttt{final\_metrics\_atlas\_tagger\_xgb\_final\_production\_round1\_round1.json};
  \item paired SHAP-Delta campaign \texttt{shap\_delta\_final\_production\_round3};
  \item conformal campaign \texttt{conformal\_coverage\_round1};
  \item pairing logs \texttt{ATLAS80030\_pair\_all\_round1} and
        \texttt{ATLAS80030\_pair\_bias\_round1}.
\end{itemize}

Tables~\ref{tab:supp-syst-detector}--\ref{tab:supp-shap-scores} are produced
from those artifacts by \texttt{tables/make\_supplementary\_tables.py}.
Table~\ref{tab:supp-stale-hlf} comes from the completed full-sample
\texttt{03\_Code\_files/stale\_hlf\_audit.py} outputs. The re-grouping script for
Tables~\ref{tab:supp-syst-detector}--\ref{tab:supp-shap-scores} copies every
number verbatim and asserts that the multiset of numeric tokens in its outputs
is exactly the multiset in the frozen sources, so a value cannot be dropped,
altered or introduced by the layout change. The paired SHAP-Delta outputs are
additionally cross-checked value by value against the frozen twenty-column
artifact of the previous revision.

\begin{description}
  \item[Table~\ref{tab:supp-stale-hlf}] the stored high-level branch audit that
    supports Section~2.3 of the main text.
  \item[Tables~\ref{tab:supp-syst-detector} and \ref{tab:supp-syst-modelling}]
    the bin-by-bin systematic breakdown underlying Fig.~2(b) of the main text,
    split into its detector part and its modelling and scale part.
  \item[Tables~\ref{tab:supp-conformal-coverage} and \ref{tab:supp-conformal-sets}]
    the complete conformal result underlying Fig.~4 of the main text, split into
    coverage and prediction-set structure.
  \item[Tables~\ref{tab:supp-shap-per-feature} and \ref{tab:supp-shap-scores}]
    the complete paired SHAP-Delta result underlying Fig.~3 and Table~5 of the
    main text, split into per-feature attribution results and systematic-level
    score and rank summaries.
\end{description}

\clearpage
\section{Stored high-level branches in the varied samples}

For the four calorimeter-cluster variations with validated exact pairing, the
fifteen stored \texttt{fjet\_*} high-level branches are float32 bit-identical to
nominal for every matched jet, although the released constituents change.
Table~\ref{tab:supp-stale-hlf} reports $5{,}000{,}000$ exact matches for each of
\texttt{esup}, \texttt{esdown} and \texttt{cpos}, and $4{,}995{,}059$ for
\texttt{cer}. Across all eligible matched-jet/branch comparisons, the exact
identity counts are $75{,}000{,}000/75{,}000{,}000$ for each fully matched
sample and $74{,}925{,}885/74{,}925{,}885$ for \texttt{cer}.

The relative summed-constituent-$\pt$ shift has mean $+6.780753\%$ and median
$+7.406764\%$ under \texttt{esup}; the corresponding \texttt{esdown} mean and
median are $-6.780749\%$ and $-7.406822\%$. The \texttt{cer} summed-energy
response has RMS $3.770123\%$, and the median \texttt{cpos}
constituent-centroid displacement is $\Delta R=4.9387\times10^{-5}$.
The track samples and \texttt{bias} are not pairable on the full sample:
their exact-key-hit counts are diagnostic only. Track constituent
multiplicity is therefore compared as an unpaired sample distribution;
no pairwise stored-branch identity or same-jet constituent-change fraction
is inferred for those five samples.

% Audited values from 04_Results/stale_hlf_audit/
% supp_s1_stale_hlf_audit.tex, with only layout and explanatory text adapted
% for the supplementary document. No physics-analysis result is recalculated.
\begin{table*}[!ht]
\centering\footnotesize
\caption{Full-sample audit of stored high-level branches in the nine detector
  variations. The four calorimeter-cluster samples pass the exact-key pairing
  preflight; for every valid matched jet, all fifteen stored high-level values
  are identical to nominal as float32 bit patterns. Exact-key hits for track
  samples and \texttt{bias} are diagnostics only and do not support pairwise
  identity fractions.}
\label{tab:supp-stale-hlf}
\setlength{\tabcolsep}{3pt}
\resizebox{\textwidth}{!}{%
\begin{tabular}{lrrlll}
\toprule
Sample & Shards & Exact-key hits & Pairing basis & HLF identity & Constituent evidence \\
\midrule
\texttt{esup} & 50 & 5{,}000{,}000 & exact key (100.000\%) & 75{,}000{,}000/75{,}000{,}000 (100.000\%) & median $\Delta\Sigma p_{\mathrm T}$: +7.4068\% \\
\texttt{esdown} & 50 & 5{,}000{,}000 & exact key (100.000\%) & 75{,}000{,}000/75{,}000{,}000 (100.000\%) & median $\Delta\Sigma p_{\mathrm T}$: -7.4068\% \\
\texttt{cer} & 50 & 4{,}995{,}059 & exact key (99.901\%) & 74{,}925{,}885/74{,}925{,}885 (100.000\%) & RMS $\Delta\Sigma E$: 3.7701\% \\
\texttt{cpos} & 50 & 5{,}000{,}000 & exact key (100.000\%) & 75{,}000{,}000/75{,}000{,}000 (100.000\%) & median centroid $\Delta R$: $4.9387\times10^{-5}$ \\
\texttt{teg} & 50 & 2{,}250{,}423 & N/A / not pairable & N/A & median $N_{\mathrm c}$: 49 vs nominal 50 (unpaired) \\
\texttt{tej} & 50 & 1{,}906{,}619 & N/A / not pairable & N/A & median $N_{\mathrm c}$: 49 vs nominal 50 (unpaired) \\
\texttt{tfl} & 50 & 2{,}749{,}571 & N/A / not pairable & N/A & median $N_{\mathrm c}$: 48 vs nominal 50 (unpaired) \\
\texttt{tfj} & 50 & 2{,}006{,}933 & N/A / not pairable & N/A & median $N_{\mathrm c}$: 49 vs nominal 50 (unpaired) \\
\texttt{bias} & 50 & 4{,}167{,}067 & N/A / not pairable & N/A & N/A \\
\bottomrule
\end{tabular}%
}
\par\smallskip
\begin{minipage}{\textwidth}\footnotesize
\emph{Definitions:} $\Delta\Sigma p_{\mathrm T}=\Sigma p_{\mathrm T}^{\rm sys}/
\Sigma p_{\mathrm T}^{\rm nom}-1$; $\Delta\Sigma E=\Sigma E^{\rm sys}/
\Sigma E^{\rm nom}-1$. The \texttt{cer} value is the RMS of the signed
summed-energy response, not its near-zero signed median. The \texttt{cpos}
value is the median $\Delta R$ of the constituent-$p_{\mathrm T}$-weighted
$\eta$--$\phi$ centroids. Track multiplicities are separate full-sample
distributions, not matched-jet changes; the nominal median is 50.
\texttt{teg} has 183{,}802 duplicate keys, and the other track samples
and \texttt{bias} fall below the validated pairing-coverage floor.
The full counts and per-branch comparisons are preserved in
\texttt{04\_Results/stale\_hlf\_audit/}.
\end{minipage}
\end{table*}


\section{Bin-by-bin systematic uncertainties}

The fourteen physical $\pt$ bins span $350$--$3150$~GeV in $200$~GeV steps. All
uncertainty columns are relative uncertainties on background rejection at $50\%$
signal efficiency, in percent. Table~\ref{tab:supp-syst-detector} carries the
cluster and track categories, Table~\ref{tab:supp-syst-modelling} the generator
modelling and ISR/FSR scale categories together with the total of all four. No
statistical uncertainty is attached to any entry.

% Generated by tables/make_supplementary_tables.py from the frozen
% artifacts. Every value is copied verbatim; see PROVENANCE.md.
% Requires \usepackage{booktabs}, pdflscape.
\begin{landscape}
\begin{table}[p]
\centering
\footnotesize
\setlength{\tabcolsep}{4pt}
\caption{Detector systematic uncertainties on background rejection at $50\%$ signal efficiency, in each of the fourteen physical $200$~GeV $\pt$ bins. All uncertainty columns are relative uncertainties in percent. The nominal background rejection is repeated here so that each relative uncertainty can be read against the value it applies to.}
\label{tab:supp-syst-detector}
\begin{tabular}{rrr rrrr rrrr}
\toprule
$\pt$ low & $\pt$ high & Nominal BR & CES & CER & CPOS & Cluster total & Track eff. & Fake rate & Bias & Track total \\
\midrule
350 & 550 & 93.414 & 9.97923 & 4.11817 & 1.54679 & 10.9058 & 3.78563 & 2.07996 & 1.73344 & 4.65425 \\
550 & 750 & 66.9929 & 11.3853 & 2.31711 & 0.489907 & 11.629 & 3.15357 & 4.43777 & 0.608945 & 5.4781 \\
750 & 950 & 58.5127 & 14.3254 & 3.8968 & 1.48534 & 14.9201 & 4.0226 & 2.93045 & 3.05909 & 5.84182 \\
950 & 1150 & 53.7758 & 15.5448 & 1.04451 & 0.0182623 & 15.5799 & 1.72603 & 2.74085 & 0.92674 & 3.36901 \\
1150 & 1350 & 55.0452 & 13.8141 & 3.44808 & 1.34744 & 14.3016 & 1.42879 & 4.61806 & 2.5554 & 5.4679 \\
1350 & 1550 & 48.4567 & 13.3623 & 2.79696 & 0.909349 & 13.6822 & 2.99679 & 4.45887 & 1.76832 & 5.6559 \\
1550 & 1750 & 49.1348 & 13.2405 & 3.47274 & 0.837172 & 13.7139 & 0.915461 & 2.82229 & 0.175584 & 2.97224 \\
1750 & 1950 & 44.1886 & 14.2629 & 4.77981 & 1.38144 & 15.1058 & 1.54494 & 2.68994 & 0.422248 & 3.13064 \\
1950 & 2150 & 40.8557 & 9.92577 & 1.8989 & 0.812175 & 10.1384 & 0.921182 & 4.53883 & 1.96547 & 5.03116 \\
2150 & 2350 & 39.6283 & 12.6348 & 3.44452 & 1.26718 & 13.1571 & 4.11698 & 6.04944 & 2.89518 & 7.8694 \\
2350 & 2550 & 37.5143 & 12.5682 & 2.31997 & 0.801577 & 12.8056 & 1.73761 & 5.46226 & 0.90764 & 5.80339 \\
2550 & 2750 & 33.712 & 6.65472 & 1.2474 & 0.558901 & 6.79365 & 0.910098 & 4.01797 & 0.113344 & 4.12131 \\
2750 & 2950 & 32.3081 & 8.18673 & 4.0166 & 1.71849 & 9.27948 & 2.19987 & 4.20073 & 2.37679 & 5.30421 \\
2950 & 3150 & 28.5641 & 7.27987 & 0.0548925 & 0.259853 & 7.28472 & 3.14807 & 4.78966 & 0.284346 & 5.73864 \\
\bottomrule
\end{tabular}
\par\vspace{4pt}
\begin{minipage}{\textwidth}\footnotesize $\pt$ bin edges in GeV. CES is the \texttt{esup}/\texttt{esdown} envelope, track efficiency the \texttt{teg}/\texttt{tej} envelope and fake rate the \texttt{tfl}/\texttt{tfj} envelope; category totals are quadrature sums of the components to their left.\end{minipage}
\end{table}
\end{landscape}


% Generated by tables/make_supplementary_tables.py from the frozen
% artifacts. Every value is copied verbatim; see PROVENANCE.md.
% Requires \usepackage{booktabs}, pdflscape.
\begin{landscape}
\begin{table}[p]
\centering
\footnotesize
\setlength{\tabcolsep}{4pt}
\caption{Generator-modelling and ISR/FSR scale uncertainties on background rejection at $50\%$ signal efficiency, in the same fourteen bins, with the total of all four categories. All columns are relative uncertainties in percent.}
\label{tab:supp-syst-modelling}
\begin{tabular}{rr rrr rrrr r r}
\toprule
$\pt$ low & $\pt$ high & Signal model & Bkg model & Model total & Signal ISR & Signal FSR & Bkg ISR & Bkg FSR & ISR/FSR total & Total \\
\midrule
350 & 550 & 20.2801 & 20.2063 & 28.6283 & 1.29504 & 20.399 & 7.22165 & 10.5631 & 24.1149 & 39.2645 \\
550 & 750 & 14.1618 & 9.56575 & 17.0898 & 0.774375 & 18.349 & 6.88518 & 7.21155 & 20.8973 & 29.8999 \\
750 & 950 & 8.39876 & 10.374 & 13.3476 & 0.528659 & 15.4347 & 5.0622 & 8.73264 & 18.4497 & 27.844 \\
950 & 1150 & 9.35958 & 6.08504 & 11.1638 & 0.352858 & 16.4144 & 0.518628 & 5.87127 & 17.4441 & 26.1344 \\
1150 & 1350 & 10.7978 & 1.46415 & 10.8966 & 0.440081 & 17.2831 & 3.44701 & 5.70108 & 18.5279 & 26.3904 \\
1350 & 1550 & 10.7978 & 7.20897 & 12.9831 & 0.289097 & 17.3928 & 1.57054 & 5.95989 & 18.4548 & 26.9876 \\
1550 & 1750 & 20.2801 & 7.25911 & 21.5401 & 0.244223 & 11.1576 & 1.89294 & 4.79582 & 12.2937 & 28.4959 \\
1750 & 1950 & 20.2801 & 5.62898 & 21.0468 & 0.21374 & 10.6371 & 1.91893 & 3.8948 & 11.4911 & 28.5132 \\
1950 & 2150 & 20.2801 & 4.21909 & 20.7143 & 0.367505 & 16.1712 & 0.924417 & 1.12223 & 16.2406 & 28.6521 \\
2150 & 2350 & 20.2801 & 4.25508 & 20.7217 & 0.149738 & 14.252 & 1.18127 & 1.59226 & 14.3901 & 29.5212 \\
2350 & 2550 & 20.2801 & 4.69118 & 20.8156 & 0.145985 & 8.02331 & 0.245851 & 1.03581 & 8.09494 & 26.3909 \\
2550 & 2750 & 20.2801 & 7.14007 & 21.5003 & 0.0935954 & 10.6968 & 2.03069 & 0.639324 & 10.907 & 25.3844 \\
2750 & 2950 & 20.2801 & 6.17343 & 21.1989 & 0.0594766 & 11.9783 & 2.16768 & 0.957883 & 12.2107 & 26.6972 \\
2950 & 3150 & 20.2801 & 6.45282 & 21.282 & 0.0653595 & 5.29897 & 1.107 & 0.457086 & 5.43302 & 23.842 \\
\bottomrule
\end{tabular}
\par\vspace{4pt}
\begin{minipage}{\textwidth}\footnotesize $\pt$ bin edges in GeV. Signal modelling is measured only in the bins below $1.5$~TeV and held at its maximum measured value from $1.55$~TeV upwards; the $1350$--$1550$~GeV interval carries the official plotting convention described in Section~4.1 of the main text and is not an independent measurement of that interval. The Total column is the quadrature sum of the cluster and track totals of Table~\ref{tab:supp-syst-detector} with the model and ISR/FSR totals here.\end{minipage}
\end{table}
\end{landscape}


\section{Class-conditional split-conformal coverage}

Thresholds were calibrated once, at target coverage $0.9$, on the frozen held-out
nominal validation split ($500{,}385$ jets; $252{,}215$ signal, $248{,}170$
background), giving $q_{\mathrm{sig}} = 0.5700125098228455$ and
$q_{\mathrm{bkg}} = 0.6447482109069824$ on the nonconformity score
$s_c = 1 - \hat p_c$. Class $c$ enters a jet's prediction set when
$1 - \hat p_c \le q_c$, equivalently $\hat p_c \ge 1 - q_c$. The thresholds were
applied unchanged to every sample below. Coverage is empirical and unweighted; no
refitting, recalibration or reweighting was performed, and the calibration jets
took no part in fitting the model or in selecting its hyperparameters.

The two class-conditional coverages are the primary quantities. Overall coverage
is a mixture of them weighted by each sample's class composition, which ranges
from $50.5\%$ signal on \texttt{test\_nominal} to $18.5\%$ on \texttt{tfl}, and is
reported as a descriptive number only. Deviations are quoted both from the $0.9$
target and from the value measured on the nominal test sample.

% Generated by tables/make_supplementary_tables.py from the frozen
% artifacts. Every value is copied verbatim; see PROVENANCE.md.
% Requires \usepackage{booktabs}, pdflscape.
\begin{landscape}
\begin{table}[p]
\centering
\footnotesize
\setlength{\tabcolsep}{3pt}
\caption{Class-conditional split-conformal coverage at target $0.9$. The two class-conditional coverages are the primary quantities; the overall coverage is a descriptive number whose value depends on the class mixture of the sample, which differs strongly between the calorimeter and the track samples.}
\label{tab:supp-conformal-coverage}
\begin{tabular}{l rrr rrr rrr rrr}
\toprule
Dataset & $n$ & $n_s$ & $n_b$ & Signal cov. & Bkg cov. & Overall cov. & $\Delta_s$ target & $\Delta_b$ target & $\Delta$ target & $\Delta_s$ nominal & $\Delta_b$ nominal & $\Delta$ nominal \\
\midrule
test\_nominal & 10306812 & 5207478 & 5099334 & 0.900423 & 0.8999877 & 0.9002076 & 0.0004230071 & -1.227611e-05 & 0.0002076491 & 0 & 0 & 0 \\
esup & 4999999 & 2592575 & 2407424 & 0.9011597 & 0.8960561 & 0.8987024 & 0.001159658 & -0.003943884 & -0.00129762 & 0.0007366508 & -0.003931607 & -0.001505269 \\
esdown & 4999999 & 2592653 & 2407346 & 0.8903914 & 0.9021154 & 0.8960362 & -0.009608575 & 0.002115442 & -0.003963821 & -0.01003158 & 0.002127718 & -0.00417147 \\
cer & 4999756 & 2590441 & 2409315 & 0.8989906 & 0.8994121 & 0.8991937 & -0.001009442 & -0.0005879264 & -0.0008063193 & -0.001432449 & -0.0005756503 & -0.001013968 \\
cpos & 4999999 & 2592597 & 2407402 & 0.9002356 & 0.8994331 & 0.8998492 & 0.0002355553 & -0.0005669182 & -0.00015082 & -0.0001874518 & -0.0005546421 & -0.0003584691 \\
teg & 4999998 & 1500292 & 3499706 & 0.8994109 & 0.8981366 & 0.898519 & -0.0005890853 & -0.001863414 & -0.001481041 & -0.001012092 & -0.001851138 & -0.00168869 \\
tej & 4999999 & 1589974 & 3410025 & 0.9004927 & 0.8999667 & 0.900134 & 0.0004927125 & -3.328421e-05 & 0.00013398 & 6.970539e-05 & -2.10081e-05 & -7.366905e-05 \\
tfl & 4999996 & 923044 & 4076952 & 0.898085 & 0.8999292 & 0.8995887 & -0.001914968 & -7.083723e-05 & -0.0004112803 & -0.002337975 & -5.856112e-05 & -0.0006189294 \\
tfj & 4999983 & 1599126 & 3400857 & 0.8957743 & 0.9003584 & 0.8988923 & -0.004225683 & 0.0003583509 & -0.001107744 & -0.00464869 & 0.000370627 & -0.001315393 \\
bias & 4999999 & 2597808 & 2402191 & 0.9005931 & 0.8997295 & 0.9001782 & 0.0005931154 & -0.0002705447 & 0.00017818 & 0.0001701084 & -0.0002582686 & -2.946904e-05 \\
\bottomrule
\end{tabular}
\par\vspace{4pt}
\begin{minipage}{\textwidth}\footnotesize Deviations are quoted as fractions, from the $0.9$ target and from the value measured on \texttt{test\_nominal}. Thresholds were calibrated once on the frozen nominal validation split and applied unchanged to every row.\end{minipage}
\end{table}
\end{landscape}


% Generated by tables/make_supplementary_tables.py from the frozen
% artifacts. Every value is copied verbatim; see PROVENANCE.md.
% Requires \usepackage{booktabs}.
\begin{table}[!ht]
\centering
\footnotesize
\setlength{\tabcolsep}{4pt}
\caption{Structure of the class-conditional conformal prediction sets on the same samples: mean set size, and the rate and count of one-class, two-class and empty sets.}
\label{tab:supp-conformal-sets}
\begin{tabular}{l r rrr rrr}
\toprule
Dataset & Mean set size & Singleton rate & Ambiguous rate & Empty rate & Singleton $n$ & Ambiguous $n$ & Empty $n$ \\
\midrule
test\_nominal & 1.103904 & 0.8960961 & 0.1039039 & 0 & 9235894 & 1070918 & 0 \\
esup & 1.106734 & 0.893266 & 0.106734 & 0 & 4466329 & 533670 & 0 \\
esdown & 1.106356 & 0.893644 & 0.106356 & 0 & 4468219 & 531780 & 0 \\
cer & 1.104816 & 0.8951841 & 0.1048159 & 0 & 4475702 & 524054 & 0 \\
cpos & 1.104307 & 0.895693 & 0.104307 & 0 & 4478464 & 521535 & 0 \\
teg & 1.100242 & 0.8997582 & 0.1002418 & 0 & 4498789 & 501209 & 0 \\
tej & 1.099836 & 0.9001636 & 0.09983642 & 0 & 4500817 & 499182 & 0 \\
tfl & 1.097283 & 0.9027171 & 0.09728288 & 0 & 4513582 & 486414 & 0 \\
tfj & 1.100723 & 0.8992775 & 0.1007225 & 0 & 4496372 & 503611 & 0 \\
bias & 1.104136 & 0.895864 & 0.104136 & 0 & 4479319 & 520680 & 0 \\
\bottomrule
\end{tabular}
\par\vspace{4pt}
\begin{minipage}{\textwidth}\footnotesize ``Ambiguous'' denotes a prediction set containing both classes. No sample produces an empty set.\end{minipage}
\end{table}


\section{Complete paired SHAP-Delta results}

SHAP values are exact TreeSHAP attributions of the frozen model in the additive
margin space of the booster; score values are in probability space.
``Mean $|\Delta$SHAP$|$'' is the mean over jets of the absolute per-jet change in
a feature's attribution, and ``relative attribution drift'' is the relative change
in that feature's global mean $|$SHAP$|$. These are different quantities and are
reported separately.

Table~\ref{tab:supp-shap-per-feature} gives the per-feature results, which are the
quantities that differ from feature to feature. Table~\ref{tab:supp-shap-scores},
which follows it, gives the score-level and rank-level summaries, which are properties of a whole
systematic, sample-size and class-subset group; in the previous revision these
were repeated on all fifteen feature rows of each group, and separating them
removes that repetition.

% Generated by tables/make_supplementary_tables.py from the frozen
% shap_delta_per_feature.csv of round 3. Values copied verbatim.
% Requires \usepackage{booktabs,longtable,pdflscape}.
\begin{landscape}
\scriptsize
\setlength{\tabcolsep}{1.6pt}
\begin{longtable}{llll r rrrr rr r rr}
\caption{Complete per-feature paired SHAP-Delta results, round 3. SHAP values are exact TreeSHAP attributions of the frozen model in the additive margin space of the booster. Mean $|\Delta$SHAP$|$ is the mean over jets of the absolute per-jet change in a feature's attribution; relative attribution drift is the relative change in that feature's global mean $|$SHAP$|$. The two are different quantities and are not interchangeable. Score-level and rank-level summaries, which are properties of a whole systematic--subset--sample-size group rather than of a feature, are given separately in Table~\ref{tab:supp-shap-scores}.}
\label{tab:supp-shap-per-feature}\\
\toprule
Systematic & Size & Subset & Feature & $n$ & Mean $\Delta$SHAP & Median $\Delta$SHAP & Mean $|\Delta$SHAP$|$ & Median $|\Delta$SHAP$|$ & Nom. mean $|$SHAP$|$ & Sys. mean $|$SHAP$|$ & Rel. attr. drift & Nom. rank & Sys. rank \\
\midrule
\endfirsthead
\toprule
Systematic & Size & Subset & Feature & $n$ & Mean $\Delta$SHAP & Median $\Delta$SHAP & Mean $|\Delta$SHAP$|$ & Median $|\Delta$SHAP$|$ & Nom. mean $|$SHAP$|$ & Sys. mean $|$SHAP$|$ & Rel. attr. drift & Nom. rank & Sys. rank \\
\midrule
\endhead
\bottomrule
\endfoot
esup & 25000 & all & $\tau_1^{\mathrm{WTA}}$ & 25000 & -0.00434156 & -0.00425814 & 0.0190665 & 0.012869 & 0.240339 & 0.235264 & -0.0211173 & 11 & 11 \\
esup & 25000 & all & $\tau_2^{\mathrm{WTA}}$ & 25000 & -0.00731675 & -0.00515284 & 0.0194375 & 0.0137851 & 0.28977 & 0.293205 & 0.0118541 & 6 & 6 \\
esup & 25000 & all & $\tau_3^{\mathrm{WTA}}$ & 25000 & -0.00367214 & -0.0020184 & 0.0223208 & 0.0146312 & 0.263671 & 0.267294 & 0.0137407 & 10 & 9 \\
esup & 25000 & all & $\tau_4^{\mathrm{WTA}}$ & 25000 & 0.0181638 & 0.0142685 & 0.0461951 & 0.0339434 & 0.915571 & 0.915845 & 0.000299139 & 1 & 1 \\
esup & 25000 & all & $\sqrt{d_{12}}$ & 25000 & 0.0109114 & 0.00742533 & 0.0353833 & 0.0249318 & 0.264725 & 0.262701 & -0.00764745 & 9 & 10 \\
esup & 25000 & all & $\sqrt{d_{23}}$ & 25000 & 0.00694671 & 0.00527251 & 0.0418826 & 0.0294513 & 0.382672 & 0.377808 & -0.0127097 & 5 & 5 \\
esup & 25000 & all & ECF$_1$ & 25000 & -0.0045179 & -0.00805226 & 0.0513568 & 0.0284781 & 0.152907 & 0.152777 & -0.000844817 & 12 & 12 \\
esup & 25000 & all & ECF$_2$ & 25000 & 0.00284624 & 0.00170189 & 0.0153989 & 0.0095173 & 0.044636 & 0.0456228 & 0.0221074 & 15 & 15 \\
esup & 25000 & all & ECF$_3$ & 25000 & -0.013611 & 0.0101904 & 0.128456 & 0.0773123 & 0.712902 & 0.677895 & -0.0491062 & 2 & 2 \\
esup & 25000 & all & $C_2$ & 25000 & 0.000575124 & 6.85288e-05 & 0.0103149 & 0.00631687 & 0.0960636 & 0.0958936 & -0.00176997 & 14 & 14 \\
esup & 25000 & all & $D_2$ & 25000 & -0.0114796 & -0.00908495 & 0.0240718 & 0.0167221 & 0.285947 & 0.280028 & -0.0207 & 7 & 7 \\
esup & 25000 & all & $L_2$ & 25000 & -0.00247824 & -0.00235188 & 0.0472834 & 0.029814 & 0.411918 & 0.428734 & 0.0408259 & 3 & 3 \\
esup & 25000 & all & $L_3$ & 25000 & 0.00307768 & 0.00332733 & 0.0204141 & 0.013195 & 0.139444 & 0.136245 & -0.0229404 & 13 & 13 \\
esup & 25000 & all & $Q_w$ & 25000 & 0.0177802 & 0.0116571 & 0.0402339 & 0.0263597 & 0.271452 & 0.272451 & 0.00368286 & 8 & 8 \\
esup & 25000 & all & ThrustMaj & 25000 & -0.0112654 & -0.00611538 & 0.0472179 & 0.0321531 & 0.409828 & 0.41075 & 0.00224971 & 4 & 4 \\
esup & 25000 & signal & $\tau_1^{\mathrm{WTA}}$ & 12961 & -0.00742859 & -0.00598618 & 0.0180438 & 0.0129641 & 0.198022 & 0.190008 & -0.0404674 & 10 & 11 \\
esup & 25000 & signal & $\tau_2^{\mathrm{WTA}}$ & 12961 & -0.00869531 & -0.00659244 & 0.0192281 & 0.0141659 & 0.260284 & 0.26242 & 0.00820569 & 8 & 8 \\
esup & 25000 & signal & $\tau_3^{\mathrm{WTA}}$ & 12961 & -0.00495204 & -0.00314064 & 0.0247753 & 0.0168893 & 0.292879 & 0.297476 & 0.0156964 & 6 & 7 \\
esup & 25000 & signal & $\tau_4^{\mathrm{WTA}}$ & 12961 & 0.020185 & 0.0155597 & 0.0469984 & 0.0351849 & 0.877359 & 0.873818 & -0.00403692 & 1 & 1 \\
esup & 25000 & signal & $\sqrt{d_{12}}$ & 12961 & 0.00941356 & 0.00743981 & 0.0334983 & 0.0241757 & 0.197561 & 0.202804 & 0.0265387 & 11 & 9 \\
esup & 25000 & signal & $\sqrt{d_{23}}$ & 12961 & 0.000580019 & -0.000811756 & 0.0446346 & 0.0326138 & 0.351147 & 0.347798 & -0.00953497 & 5 & 5 \\
esup & 25000 & signal & ECF$_1$ & 12961 & -0.0150534 & -0.0115825 & 0.048538 & 0.0329386 & 0.146334 & 0.140155 & -0.0422269 & 12 & 12 \\
esup & 25000 & signal & ECF$_2$ & 12961 & 0.00137556 & 0.000689186 & 0.0138125 & 0.00890921 & 0.0342044 & 0.0359198 & 0.0501518 & 15 & 15 \\
esup & 25000 & signal & ECF$_3$ & 12961 & -0.054434 & -0.0200152 & 0.164731 & 0.102863 & 0.628236 & 0.585732 & -0.0676562 & 2 & 2 \\
esup & 25000 & signal & $C_2$ & 12961 & 0.000765133 & 0.000294112 & 0.00950268 & 0.00608435 & 0.0804367 & 0.0800343 & -0.00500249 & 14 & 13 \\
esup & 25000 & signal & $D_2$ & 12961 & -0.011811 & -0.0107869 & 0.0230544 & 0.0171047 & 0.20712 & 0.193486 & -0.0658275 & 9 & 10 \\
esup & 25000 & signal & $L_2$ & 12961 & 0.00349563 & -0.00143513 & 0.057718 & 0.0377151 & 0.407798 & 0.428608 & 0.0510304 & 4 & 4 \\
esup & 25000 & signal & $L_3$ & 12961 & 0.00360274 & 0.00409439 & 0.015993 & 0.0109154 & 0.0821967 & 0.0772485 & -0.0602001 & 13 & 14 \\
esup & 25000 & signal & $Q_w$ & 12961 & 0.0184303 & 0.0123477 & 0.0489395 & 0.0353009 & 0.288984 & 0.297729 & 0.0302614 & 7 & 6 \\
esup & 25000 & signal & ThrustMaj & 12961 & -0.0177852 & -0.0114986 & 0.0565996 & 0.0416927 & 0.446352 & 0.441174 & -0.0116006 & 3 & 3 \\
esup & 25000 & background & $\tau_1^{\mathrm{WTA}}$ & 12039 & -0.00101812 & -0.00237626 & 0.0201676 & 0.0127597 & 0.285895 & 0.283984 & -0.00668402 & 9 & 9 \\
esup & 25000 & background & $\tau_2^{\mathrm{WTA}}$ & 12039 & -0.00583262 & -0.00374976 & 0.019663 & 0.0132174 & 0.321513 & 0.326348 & 0.0150382 & 8 & 8 \\
esup & 25000 & background & $\tau_3^{\mathrm{WTA}}$ & 12039 & -0.00229422 & -0.00119336 & 0.0196784 & 0.0125813 & 0.232227 & 0.234798 & 0.0110714 & 11 & 11 \\
esup & 25000 & background & $\tau_4^{\mathrm{WTA}}$ & 12039 & 0.0159877 & 0.0132152 & 0.0453303 & 0.0325543 & 0.956712 & 0.961086 & 0.00457163 & 1 & 1 \\
esup & 25000 & background & $\sqrt{d_{12}}$ & 12039 & 0.012524 & 0.0074071 & 0.0374127 & 0.0258643 & 0.337036 & 0.327184 & -0.0292321 & 7 & 7 \\
esup & 25000 & background & $\sqrt{d_{23}}$ & 12039 & 0.013801 & 0.0101193 & 0.0389199 & 0.0259463 & 0.416606 & 0.410117 & -0.0155771 & 3 & 4 \\
esup & 25000 & background & ECF$_1$ & 12039 & 0.00682441 & -0.00510362 & 0.0543914 & 0.0241138 & 0.159984 & 0.166366 & 0.0398924 & 13 & 13 \\
esup & 25000 & background & ECF$_2$ & 12039 & 0.00442955 & 0.00289722 & 0.0171068 & 0.0103102 & 0.0558665 & 0.0560688 & 0.00362197 & 15 & 15 \\
esup & 25000 & background & ECF$_3$ & 12039 & 0.0303384 & 0.0319794 & 0.0894031 & 0.0613401 & 0.804059 & 0.777115 & -0.0335097 & 2 & 2 \\
esup & 25000 & background & $C_2$ & 12039 & 0.000370563 & -0.000122786 & 0.0111893 & 0.0066075 & 0.112887 & 0.112969 & 0.000721118 & 14 & 14 \\
esup & 25000 & background & $D_2$ & 12039 & -0.0111228 & -0.00714368 & 0.0251672 & 0.0163062 & 0.37081 & 0.373197 & 0.00643682 & 5 & 6 \\
esup & 25000 & background & $L_2$ & 12039 & -0.00890962 & -0.002859 & 0.0360496 & 0.0235629 & 0.416353 & 0.428868 & 0.0300589 & 4 & 3 \\
esup & 25000 & background & $L_3$ & 12039 & 0.00251241 & 0.00220869 & 0.0251738 & 0.0168725 & 0.201075 & 0.19976 & -0.00653881 & 12 & 12 \\
esup & 25000 & background & $Q_w$ & 12039 & 0.0170803 & 0.0111133 & 0.0308616 & 0.0192962 & 0.252578 & 0.245239 & -0.0290558 & 10 & 10 \\
esup & 25000 & background & ThrustMaj & 12039 & -0.00424623 & -0.0027098 & 0.0371177 & 0.0243137 & 0.370504 & 0.377996 & 0.0202232 & 6 & 5 \\
esup & 50000 & all & $\tau_1^{\mathrm{WTA}}$ & 50000 & -0.00435308 & -0.00434025 & 0.0191103 & 0.0129801 & 0.239952 & 0.234746 & -0.0216935 & 11 & 11 \\
esup & 50000 & all & $\tau_2^{\mathrm{WTA}}$ & 50000 & -0.00739281 & -0.00518003 & 0.0193923 & 0.0136991 & 0.289399 & 0.29289 & 0.0120618 & 6 & 6 \\
esup & 50000 & all & $\tau_3^{\mathrm{WTA}}$ & 50000 & -0.00391045 & -0.00221208 & 0.0224239 & 0.0148252 & 0.265045 & 0.26864 & 0.0135639 & 9 & 9 \\
esup & 50000 & all & $\tau_4^{\mathrm{WTA}}$ & 50000 & 0.0178074 & 0.0143479 & 0.0462972 & 0.0338561 & 0.914396 & 0.915141 & 0.000815591 & 1 & 1 \\
esup & 50000 & all & $\sqrt{d_{12}}$ & 50000 & 0.011045 & 0.00767586 & 0.0355069 & 0.0250531 & 0.263668 & 0.261946 & -0.00653164 & 10 & 10 \\
esup & 50000 & all & $\sqrt{d_{23}}$ & 50000 & 0.0069583 & 0.00545232 & 0.042014 & 0.02936 & 0.382003 & 0.377171 & -0.0126507 & 5 & 5 \\
esup & 50000 & all & ECF$_1$ & 50000 & -0.00481076 & -0.00800223 & 0.0514146 & 0.028577 & 0.153339 & 0.153231 & -0.000702014 & 12 & 12 \\
esup & 50000 & all & ECF$_2$ & 50000 & 0.00252231 & 0.00163595 & 0.0153134 & 0.00937944 & 0.0444037 & 0.0453558 & 0.0214414 & 15 & 15 \\
esup & 50000 & all & ECF$_3$ & 50000 & -0.0128796 & 0.00989418 & 0.127355 & 0.0762991 & 0.707832 & 0.673906 & -0.0479295 & 2 & 2 \\
esup & 50000 & all & $C_2$ & 50000 & 0.000523686 & 5.70007e-05 & 0.0103817 & 0.0063462 & 0.0960296 & 0.0958756 & -0.00160371 & 14 & 14 \\
esup & 50000 & all & $D_2$ & 50000 & -0.0115232 & -0.00902383 & 0.0240372 & 0.016671 & 0.286145 & 0.280275 & -0.0205135 & 7 & 7 \\
esup & 50000 & all & $L_2$ & 50000 & -0.00247689 & -0.00230951 & 0.0471149 & 0.0299817 & 0.412361 & 0.428791 & 0.0398445 & 3 & 3 \\
esup & 50000 & all & $L_3$ & 50000 & 0.00322823 & 0.0033555 & 0.0202861 & 0.0131214 & 0.137862 & 0.134575 & -0.023845 & 13 & 13 \\
esup & 50000 & all & $Q_w$ & 50000 & 0.0176774 & 0.0118103 & 0.0402834 & 0.0266111 & 0.271573 & 0.272342 & 0.00283227 & 8 & 8 \\
esup & 50000 & all & ThrustMaj & 50000 & -0.0117808 & -0.00658655 & 0.0471487 & 0.0323781 & 0.410352 & 0.411292 & 0.00229136 & 4 & 4 \\
esup & 50000 & signal & $\tau_1^{\mathrm{WTA}}$ & 25927 & -0.00735453 & -0.00609876 & 0.0181184 & 0.0131077 & 0.19842 & 0.190313 & -0.0408576 & 10 & 11 \\
esup & 50000 & signal & $\tau_2^{\mathrm{WTA}}$ & 25927 & -0.00858766 & -0.00646585 & 0.0190256 & 0.0140086 & 0.259692 & 0.261843 & 0.00828144 & 8 & 8 \\
esup & 50000 & signal & $\tau_3^{\mathrm{WTA}}$ & 25927 & -0.00498602 & -0.00340959 & 0.0247696 & 0.0169639 & 0.294227 & 0.299049 & 0.0163885 & 6 & 6 \\
esup & 50000 & signal & $\tau_4^{\mathrm{WTA}}$ & 25927 & 0.0201126 & 0.0157146 & 0.0470576 & 0.035239 & 0.875636 & 0.872878 & -0.00314865 & 1 & 1 \\
esup & 50000 & signal & $\sqrt{d_{12}}$ & 25927 & 0.00967909 & 0.00778377 & 0.0339121 & 0.0244445 & 0.197712 & 0.20323 & 0.0279083 & 11 & 9 \\
esup & 50000 & signal & $\sqrt{d_{23}}$ & 25927 & 0.00106707 & -0.000241786 & 0.0446774 & 0.0323065 & 0.351247 & 0.34773 & -0.0100139 & 5 & 5 \\
esup & 50000 & signal & ECF$_1$ & 25927 & -0.0154546 & -0.0118567 & 0.0486477 & 0.0330138 & 0.1458 & 0.139754 & -0.0414668 & 12 & 12 \\
esup & 50000 & signal & ECF$_2$ & 25927 & 0.0011719 & 0.000628404 & 0.013604 & 0.0087861 & 0.0342838 & 0.0357997 & 0.0442178 & 15 & 15 \\
esup & 50000 & signal & ECF$_3$ & 25927 & -0.052978 & -0.0193191 & 0.16327 & 0.102191 & 0.626673 & 0.585182 & -0.066208 & 2 & 2 \\
esup & 50000 & signal & $C_2$ & 25927 & 0.000799795 & 0.00027226 & 0.00956625 & 0.0061269 & 0.0802279 & 0.0798414 & -0.00481686 & 14 & 13 \\
esup & 50000 & signal & $D_2$ & 25927 & -0.011702 & -0.0106751 & 0.0229883 & 0.0170517 & 0.207092 & 0.193448 & -0.0658812 & 9 & 10 \\
esup & 50000 & signal & $L_2$ & 25927 & 0.00345814 & -0.00119814 & 0.0573852 & 0.0378829 & 0.407063 & 0.427381 & 0.0499142 & 4 & 4 \\
esup & 50000 & signal & $L_3$ & 25927 & 0.00369082 & 0.00406937 & 0.0160022 & 0.0108565 & 0.0821137 & 0.077204 & -0.0597914 & 13 & 14 \\
esup & 50000 & signal & $Q_w$ & 25927 & 0.0179536 & 0.0126648 & 0.0488631 & 0.0352651 & 0.289722 & 0.2979 & 0.0282283 & 7 & 7 \\
esup & 50000 & signal & ThrustMaj & 25927 & -0.0184363 & -0.0123952 & 0.0561168 & 0.0412853 & 0.447112 & 0.441616 & -0.0122928 & 3 & 3 \\
esup & 50000 & background & $\tau_1^{\mathrm{WTA}}$ & 24073 & -0.00112047 & -0.00249973 & 0.0201787 & 0.0128053 & 0.28468 & 0.282601 & -0.0073012 & 9 & 9 \\
esup & 50000 & background & $\tau_2^{\mathrm{WTA}}$ & 24073 & -0.00610593 & -0.00387821 & 0.0197872 & 0.0132971 & 0.321393 & 0.326328 & 0.0153547 & 8 & 7 \\
esup & 50000 & background & $\tau_3^{\mathrm{WTA}}$ & 24073 & -0.00275205 & -0.00137648 & 0.0198975 & 0.0127311 & 0.233614 & 0.235886 & 0.0097249 & 11 & 11 \\
esup & 50000 & background & $\tau_4^{\mathrm{WTA}}$ & 24073 & 0.0153247 & 0.0130663 & 0.0454782 & 0.0323426 & 0.956137 & 0.960661 & 0.00473216 & 1 & 1 \\
esup & 50000 & background & $\sqrt{d_{12}}$ & 24073 & 0.0125161 & 0.00757877 & 0.0372245 & 0.0257534 & 0.334703 & 0.325188 & -0.0284281 & 7 & 8 \\
esup & 50000 & background & $\sqrt{d_{23}}$ & 24073 & 0.0133032 & 0.0102022 & 0.0391454 & 0.02626 & 0.415126 & 0.408883 & -0.0150386 & 4 & 4 \\
esup & 50000 & background & ECF$_1$ & 24073 & 0.00665281 & -0.0050611 & 0.0543947 & 0.0240488 & 0.161456 & 0.167748 & 0.038965 & 13 & 13 \\
esup & 50000 & background & ECF$_2$ & 24073 & 0.00397672 & 0.00283502 & 0.0171544 & 0.0101886 & 0.0553029 & 0.0556475 & 0.00623047 & 15 & 15 \\
esup & 50000 & background & ECF$_3$ & 24073 & 0.0303069 & 0.0313842 & 0.0886737 & 0.0603171 & 0.795229 & 0.769456 & -0.0324098 & 2 & 2 \\
esup & 50000 & background & $C_2$ & 24073 & 0.000226312 & -0.00012821 & 0.01126 & 0.00660009 & 0.113048 & 0.113144 & 0.000856654 & 14 & 14 \\
esup & 50000 & background & $D_2$ & 24073 & -0.0113306 & -0.00706074 & 0.025167 & 0.0162562 & 0.371285 & 0.373788 & 0.00674132 & 5 & 6 \\
esup & 50000 & background & $L_2$ & 24073 & -0.00886901 & -0.00300932 & 0.0360537 & 0.0236436 & 0.418071 & 0.430315 & 0.0292866 & 3 & 3 \\
esup & 50000 & background & $L_3$ & 24073 & 0.00273001 & 0.00229591 & 0.0249 & 0.0166543 & 0.197905 & 0.196365 & -0.0077807 & 12 & 12 \\
esup & 50000 & background & $Q_w$ & 24073 & 0.0173799 & 0.0112473 & 0.031043 & 0.0196894 & 0.252024 & 0.244814 & -0.0286081 & 10 & 10 \\
esup & 50000 & background & ThrustMaj & 24073 & -0.00461278 & -0.00284562 & 0.0374899 & 0.0247984 & 0.370759 & 0.378633 & 0.0212371 & 6 & 5 \\
esup & 100000 & all & $\tau_1^{\mathrm{WTA}}$ & 100000 & -0.00443192 & -0.00434025 & 0.0191639 & 0.0130233 & 0.239131 & 0.234035 & -0.0213118 & 11 & 11 \\
esup & 100000 & all & $\tau_2^{\mathrm{WTA}}$ & 100000 & -0.00735514 & -0.00517795 & 0.0194206 & 0.0137621 & 0.289774 & 0.293169 & 0.0117184 & 6 & 6 \\
esup & 100000 & all & $\tau_3^{\mathrm{WTA}}$ & 100000 & -0.00390298 & -0.00215878 & 0.022488 & 0.0147564 & 0.265777 & 0.269441 & 0.0137863 & 9 & 9 \\
esup & 100000 & all & $\tau_4^{\mathrm{WTA}}$ & 100000 & 0.0179069 & 0.0144962 & 0.0465714 & 0.0340235 & 0.914908 & 0.915549 & 0.000700995 & 1 & 1 \\
esup & 100000 & all & $\sqrt{d_{12}}$ & 100000 & 0.0110216 & 0.00763958 & 0.0355621 & 0.0251217 & 0.263566 & 0.261896 & -0.00633653 & 10 & 10 \\
esup & 100000 & all & $\sqrt{d_{23}}$ & 100000 & 0.00683049 & 0.00525184 & 0.0421442 & 0.0294514 & 0.381965 & 0.377015 & -0.012958 & 5 & 5 \\
esup & 100000 & all & ECF$_1$ & 100000 & -0.00520791 & -0.00800674 & 0.0512201 & 0.0283191 & 0.153874 & 0.153756 & -0.000762907 & 12 & 12 \\
esup & 100000 & all & ECF$_2$ & 100000 & 0.00248428 & 0.00159834 & 0.015399 & 0.00939831 & 0.0442762 & 0.0452584 & 0.0221844 & 15 & 15 \\
esup & 100000 & all & ECF$_3$ & 100000 & -0.0128104 & 0.00986212 & 0.127205 & 0.0756971 & 0.706964 & 0.673663 & -0.0471041 & 2 & 2 \\
esup & 100000 & all & $C_2$ & 100000 & 0.000478312 & 6.12177e-05 & 0.0104347 & 0.00640296 & 0.0960899 & 0.0959988 & -0.000948362 & 14 & 14 \\
esup & 100000 & all & $D_2$ & 100000 & -0.0115426 & -0.00901648 & 0.0239895 & 0.0166137 & 0.286172 & 0.28033 & -0.0204134 & 7 & 7 \\
esup & 100000 & all & $L_2$ & 100000 & -0.00248381 & -0.00238393 & 0.0471862 & 0.0299634 & 0.413948 & 0.430607 & 0.0402431 & 3 & 3 \\
esup & 100000 & all & $L_3$ & 100000 & 0.0032894 & 0.00334955 & 0.0202784 & 0.0130793 & 0.138227 & 0.134892 & -0.0241263 & 13 & 13 \\
esup & 100000 & all & $Q_w$ & 100000 & 0.0177597 & 0.0119417 & 0.0402534 & 0.0265733 & 0.271797 & 0.272716 & 0.00338268 & 8 & 8 \\
esup & 100000 & all & ThrustMaj & 100000 & -0.0121745 & -0.00685555 & 0.0470037 & 0.0323694 & 0.411536 & 0.412402 & 0.00210611 & 4 & 4 \\
esup & 100000 & signal & $\tau_1^{\mathrm{WTA}}$ & 51851 & -0.00739561 & -0.00601216 & 0.0181823 & 0.013091 & 0.198255 & 0.190244 & -0.0404051 & 10 & 11 \\
esup & 100000 & signal & $\tau_2^{\mathrm{WTA}}$ & 51851 & -0.00865511 & -0.00655799 & 0.0190788 & 0.014121 & 0.26058 & 0.262693 & 0.00810844 & 8 & 8 \\
esup & 100000 & signal & $\tau_3^{\mathrm{WTA}}$ & 51851 & -0.00494974 & -0.00327975 & 0.0250134 & 0.0170527 & 0.294808 & 0.299622 & 0.0163274 & 6 & 6 \\
esup & 100000 & signal & $\tau_4^{\mathrm{WTA}}$ & 51851 & 0.0202038 & 0.0158541 & 0.0475725 & 0.0354851 & 0.879015 & 0.875911 & -0.00353099 & 1 & 1 \\
esup & 100000 & signal & $\sqrt{d_{12}}$ & 51851 & 0.00964228 & 0.00775301 & 0.0339011 & 0.0244609 & 0.197514 & 0.203156 & 0.0285645 & 11 & 9 \\
esup & 100000 & signal & $\sqrt{d_{23}}$ & 51851 & 0.000961959 & -0.000287414 & 0.0449685 & 0.0325209 & 0.351549 & 0.347763 & -0.0107689 & 5 & 5 \\
esup & 100000 & signal & ECF$_1$ & 51851 & -0.0157303 & -0.0118219 & 0.0487324 & 0.0328522 & 0.145992 & 0.139981 & -0.0411721 & 12 & 12 \\
esup & 100000 & signal & ECF$_2$ & 51851 & 0.00114394 & 0.000539167 & 0.0136528 & 0.00874949 & 0.0341046 & 0.0357019 & 0.0468342 & 15 & 15 \\
esup & 100000 & signal & ECF$_3$ & 51851 & -0.0525111 & -0.0184747 & 0.16292 & 0.100925 & 0.624806 & 0.584159 & -0.0650556 & 2 & 2 \\
esup & 100000 & signal & $C_2$ & 51851 & 0.000771914 & 0.00029704 & 0.00956622 & 0.00615335 & 0.0800074 & 0.079708 & -0.00374106 & 14 & 13 \\
esup & 100000 & signal & $D_2$ & 51851 & -0.0117442 & -0.0106751 & 0.0229985 & 0.0169925 & 0.207536 & 0.193872 & -0.0658407 & 9 & 10 \\
esup & 100000 & signal & $L_2$ & 51851 & 0.00326186 & -0.00109395 & 0.0575073 & 0.0380129 & 0.407436 & 0.428088 & 0.0506868 & 4 & 4 \\
esup & 100000 & signal & $L_3$ & 51851 & 0.00362828 & 0.00397675 & 0.0160768 & 0.0108812 & 0.0818522 & 0.0769241 & -0.060207 & 13 & 14 \\
esup & 100000 & signal & $Q_w$ & 51851 & 0.0182386 & 0.0130228 & 0.0488023 & 0.0351582 & 0.289932 & 0.298148 & 0.0283356 & 7 & 7 \\
esup & 100000 & signal & ThrustMaj & 51851 & -0.018611 & -0.0126655 & 0.055957 & 0.0412667 & 0.44829 & 0.442619 & -0.0126524 & 3 & 3 \\
esup & 100000 & background & $\tau_1^{\mathrm{WTA}}$ & 48149 & -0.00124036 & -0.00262269 & 0.020221 & 0.012943 & 0.283157 & 0.281189 & -0.00694705 & 9 & 9 \\
esup & 100000 & background & $\tau_2^{\mathrm{WTA}}$ & 48149 & -0.00595522 & -0.0037545 & 0.0197887 & 0.0133459 & 0.32121 & 0.325986 & 0.0148683 & 8 & 7 \\
esup & 100000 & background & $\tau_3^{\mathrm{WTA}}$ & 48149 & -0.00277574 & -0.0013648 & 0.0197684 & 0.0125363 & 0.234514 & 0.236938 & 0.0103368 & 11 & 11 \\
esup & 100000 & background & $\tau_4^{\mathrm{WTA}}$ & 48149 & 0.0154334 & 0.0131942 & 0.0454933 & 0.0324218 & 0.953555 & 0.958226 & 0.00489786 & 1 & 1 \\
esup & 100000 & background & $\sqrt{d_{12}}$ & 48149 & 0.0125071 & 0.00753541 & 0.0373509 & 0.0259081 & 0.334694 & 0.325151 & -0.0285126 & 7 & 8 \\
esup & 100000 & background & $\sqrt{d_{23}}$ & 48149 & 0.0131502 & 0.00999898 & 0.0391028 & 0.0263354 & 0.414708 & 0.408511 & -0.0149415 & 4 & 4 \\
esup & 100000 & background & ECF$_1$ & 48149 & 0.00612354 & -0.00507588 & 0.0538991 & 0.0237906 & 0.162363 & 0.168589 & 0.0383491 & 13 & 13 \\
esup & 100000 & background & ECF$_2$ & 48149 & 0.00392767 & 0.00282866 & 0.0172796 & 0.0102395 & 0.0552302 & 0.0555497 & 0.00578487 & 15 & 15 \\
esup & 100000 & background & ECF$_3$ & 48149 & 0.0299427 & 0.0313201 & 0.088744 & 0.0603766 & 0.795439 & 0.770049 & -0.0319192 & 2 & 2 \\
esup & 100000 & background & $C_2$ & 48149 & 0.000162138 & -0.000149064 & 0.0113699 & 0.00669482 & 0.113411 & 0.113542 & 0.00115985 & 14 & 14 \\
esup & 100000 & background & $D_2$ & 48149 & -0.0113255 & -0.00710131 & 0.0250567 & 0.0161715 & 0.370859 & 0.373433 & 0.00693991 & 6 & 6 \\
esup & 100000 & background & $L_2$ & 48149 & -0.00867125 & -0.00318803 & 0.0360716 & 0.0236498 & 0.420962 & 0.433315 & 0.029346 & 3 & 3 \\
esup & 100000 & background & $L_3$ & 48149 & 0.00292446 & 0.00244758 & 0.024803 & 0.0165465 & 0.198937 & 0.197319 & -0.00813306 & 12 & 12 \\
esup & 100000 & background & $Q_w$ & 48149 & 0.017244 & 0.0113006 & 0.0310472 & 0.0197037 & 0.252271 & 0.245326 & -0.0275308 & 10 & 10 \\
esup & 100000 & background & ThrustMaj & 48149 & -0.00524318 & -0.003167 & 0.037362 & 0.024838 & 0.37196 & 0.379871 & 0.0212695 & 5 & 5 \\
esdown & 25000 & all & $\tau_1^{\mathrm{WTA}}$ & 25000 & 0.00548983 & 0.00495812 & 0.0198012 & 0.0135917 & 0.239349 & 0.244213 & 0.0203218 & 11 & 11 \\
esdown & 25000 & all & $\tau_2^{\mathrm{WTA}}$ & 25000 & 0.00718007 & 0.00518979 & 0.0203258 & 0.0143468 & 0.288449 & 0.28478 & -0.0127206 & 6 & 7 \\
esdown & 25000 & all & $\tau_3^{\mathrm{WTA}}$ & 25000 & 0.00367754 & 0.00248356 & 0.0225276 & 0.014763 & 0.265475 & 0.26222 & -0.0122599 & 9 & 10 \\
esdown & 25000 & all & $\tau_4^{\mathrm{WTA}}$ & 25000 & -0.0225386 & -0.0194148 & 0.0495251 & 0.0362575 & 0.912488 & 0.908905 & -0.00392586 & 1 & 1 \\
esdown & 25000 & all & $\sqrt{d_{12}}$ & 25000 & -0.0153461 & -0.0112449 & 0.0386302 & 0.0277124 & 0.264725 & 0.264824 & 0.000374887 & 10 & 9 \\
esdown & 25000 & all & $\sqrt{d_{23}}$ & 25000 & -0.0149912 & -0.0100492 & 0.0444465 & 0.0308905 & 0.381392 & 0.380984 & -0.00107092 & 5 & 5 \\
esdown & 25000 & all & ECF$_1$ & 25000 & 0.0098823 & 0.0106271 & 0.0504551 & 0.0293548 & 0.152646 & 0.156224 & 0.0234395 & 12 & 12 \\
esdown & 25000 & all & ECF$_2$ & 25000 & -0.00122314 & -0.00140325 & 0.0163449 & 0.00985726 & 0.0444247 & 0.0448118 & 0.00871375 & 15 & 15 \\
esdown & 25000 & all & ECF$_3$ & 25000 & -0.0327011 & -0.0263252 & 0.132458 & 0.0808765 & 0.710142 & 0.717062 & 0.00974425 & 2 & 2 \\
esdown & 25000 & all & $C_2$ & 25000 & -0.000137342 & 0.000109822 & 0.0108268 & 0.00670538 & 0.0961542 & 0.0964267 & 0.00283412 & 14 & 14 \\
esdown & 25000 & all & $D_2$ & 25000 & 0.0116923 & 0.0093205 & 0.0247355 & 0.0173449 & 0.287497 & 0.29291 & 0.0188277 & 7 & 6 \\
esdown & 25000 & all & $L_2$ & 25000 & -0.00410622 & -0.000422649 & 0.0478886 & 0.0305633 & 0.410395 & 0.389261 & -0.0514981 & 3 & 4 \\
esdown & 25000 & all & $L_3$ & 25000 & -0.00527045 & -0.00475534 & 0.0220958 & 0.0146313 & 0.139372 & 0.14392 & 0.0326331 & 13 & 13 \\
esdown & 25000 & all & $Q_w$ & 25000 & -0.0265144 & -0.0169044 & 0.0440537 & 0.0286958 & 0.272915 & 0.266294 & -0.0242626 & 8 & 8 \\
esdown & 25000 & all & ThrustMaj & 25000 & 0.00446142 & 0.00264595 & 0.0465374 & 0.0329176 & 0.409908 & 0.402456 & -0.0181783 & 4 & 3 \\
esdown & 25000 & signal & $\tau_1^{\mathrm{WTA}}$ & 12963 & 0.00776723 & 0.00669727 & 0.0185795 & 0.0134562 & 0.199079 & 0.206592 & 0.0377382 & 10 & 10 \\
esdown & 25000 & signal & $\tau_2^{\mathrm{WTA}}$ & 12963 & 0.00807515 & 0.00638862 & 0.0198695 & 0.0148277 & 0.260017 & 0.257279 & -0.0105324 & 8 & 8 \\
esdown & 25000 & signal & $\tau_3^{\mathrm{WTA}}$ & 12963 & 0.00363772 & 0.002902 & 0.024616 & 0.0168098 & 0.295426 & 0.291406 & -0.0136085 & 6 & 6 \\
esdown & 25000 & signal & $\tau_4^{\mathrm{WTA}}$ & 12963 & -0.0282572 & -0.0242871 & 0.0512555 & 0.0386373 & 0.87599 & 0.873888 & -0.00239871 & 1 & 1 \\
esdown & 25000 & signal & $\sqrt{d_{12}}$ & 12963 & -0.0171142 & -0.0138306 & 0.0383262 & 0.0282646 & 0.197745 & 0.188206 & -0.0482361 & 11 & 11 \\
esdown & 25000 & signal & $\sqrt{d_{23}}$ & 12963 & -0.0156298 & -0.00965121 & 0.0481993 & 0.0347033 & 0.349892 & 0.342878 & -0.0200452 & 5 & 5 \\
esdown & 25000 & signal & ECF$_1$ & 12963 & 0.0170601 & 0.0133149 & 0.0490129 & 0.0331023 & 0.144653 & 0.154221 & 0.0661485 & 12 & 12 \\
esdown & 25000 & signal & ECF$_2$ & 12963 & -0.000569067 & -0.000360582 & 0.013605 & 0.00912765 & 0.034188 & 0.0336373 & -0.0161079 & 15 & 15 \\
esdown & 25000 & signal & ECF$_3$ & 12963 & -0.031625 & -0.0140783 & 0.167195 & 0.108344 & 0.625003 & 0.611539 & -0.0215431 & 2 & 2 \\
esdown & 25000 & signal & $C_2$ & 12963 & -0.000321605 & -3.14713e-05 & 0.00977702 & 0.00646031 & 0.0797743 & 0.0803169 & 0.00680183 & 14 & 14 \\
esdown & 25000 & signal & $D_2$ & 12963 & 0.0112653 & 0.0105805 & 0.0232269 & 0.0175788 & 0.20827 & 0.221039 & 0.0613105 & 9 & 9 \\
esdown & 25000 & signal & $L_2$ & 12963 & -0.0163431 & -0.00645357 & 0.0578943 & 0.0384159 & 0.403392 & 0.374854 & -0.0707449 & 4 & 4 \\
esdown & 25000 & signal & $L_3$ & 12963 & -0.00636042 & -0.00581826 & 0.01851 & 0.012689 & 0.0822223 & 0.0894439 & 0.0878305 & 13 & 13 \\
esdown & 25000 & signal & $Q_w$ & 12963 & -0.0347208 & -0.0266951 & 0.0557643 & 0.0407038 & 0.289396 & 0.269266 & -0.0695591 & 7 & 7 \\
esdown & 25000 & signal & ThrustMaj & 12963 & 0.00339485 & 0.00108105 & 0.054266 & 0.0414612 & 0.446363 & 0.438703 & -0.0171612 & 3 & 3 \\
esdown & 25000 & background & $\tau_1^{\mathrm{WTA}}$ & 12037 & 0.00303724 & 0.00299107 & 0.0211169 & 0.0137612 & 0.282715 & 0.284725 & 0.00711075 & 9 & 9 \\
esdown & 25000 & background & $\tau_2^{\mathrm{WTA}}$ & 12037 & 0.00621613 & 0.00400531 & 0.0208172 & 0.0139484 & 0.319067 & 0.314399 & -0.0146316 & 8 & 8 \\
esdown & 25000 & background & $\tau_3^{\mathrm{WTA}}$ & 12037 & 0.00372044 & 0.00222018 & 0.0202785 & 0.0126773 & 0.233219 & 0.230788 & -0.0104236 & 11 & 11 \\
esdown & 25000 & background & $\tau_4^{\mathrm{WTA}}$ & 12037 & -0.01638 & -0.0146014 & 0.0476617 & 0.0336933 & 0.951788 & 0.946616 & -0.00543443 & 1 & 1 \\
esdown & 25000 & background & $\sqrt{d_{12}}$ & 12037 & -0.0134421 & -0.00805992 & 0.0389576 & 0.0271171 & 0.336857 & 0.347334 & 0.0311021 & 7 & 7 \\
esdown & 25000 & background & $\sqrt{d_{23}}$ & 12037 & -0.0143034 & -0.0105125 & 0.0404051 & 0.0273865 & 0.415312 & 0.422024 & 0.0161607 & 4 & 3 \\
esdown & 25000 & background & ECF$_1$ & 12037 & 0.00215228 & 0.00806874 & 0.0520082 & 0.0258398 & 0.161256 & 0.15838 & -0.0178372 & 13 & 13 \\
esdown & 25000 & background & ECF$_2$ & 12037 & -0.00192753 & -0.0026983 & 0.0192955 & 0.0108105 & 0.0554486 & 0.0568459 & 0.025199 & 15 & 15 \\
esdown & 25000 & background & ECF$_3$ & 12037 & -0.03386 & -0.035546 & 0.0950476 & 0.0647376 & 0.801836 & 0.830702 & 0.0360006 & 2 & 2 \\
esdown & 25000 & background & $C_2$ & 12037 & 6.10956e-05 & 0.000286982 & 0.0119574 & 0.00701508 & 0.113794 & 0.113776 & -0.000153407 & 14 & 14 \\
esdown & 25000 & background & $D_2$ & 12037 & 0.0121521 & 0.0077296 & 0.0263602 & 0.0170687 & 0.37282 & 0.37031 & -0.00673322 & 5 & 5 \\
esdown & 25000 & background & $L_2$ & 12037 & 0.00907205 & 0.00378862 & 0.0371132 & 0.0241248 & 0.417938 & 0.404776 & -0.0314908 & 3 & 4 \\
esdown & 25000 & background & $L_3$ & 12037 & -0.00409662 & -0.00309901 & 0.0259574 & 0.0172921 & 0.200919 & 0.202587 & 0.0083036 & 12 & 12 \\
esdown & 25000 & background & $Q_w$ & 12037 & -0.0176768 & -0.0115598 & 0.0314422 & 0.0199651 & 0.255166 & 0.263095 & 0.0310741 & 10 & 10 \\
esdown & 25000 & background & ThrustMaj & 12037 & 0.00561005 & 0.00348935 & 0.0382142 & 0.0253082 & 0.370652 & 0.36342 & -0.0195111 & 6 & 6 \\
esdown & 50000 & all & $\tau_1^{\mathrm{WTA}}$ & 50000 & 0.00544915 & 0.00497052 & 0.0198721 & 0.0135859 & 0.23978 & 0.244832 & 0.0210682 & 11 & 11 \\
esdown & 50000 & all & $\tau_2^{\mathrm{WTA}}$ & 50000 & 0.00730408 & 0.0051336 & 0.0202522 & 0.0144204 & 0.289092 & 0.28537 & -0.0128729 & 6 & 7 \\
esdown & 50000 & all & $\tau_3^{\mathrm{WTA}}$ & 50000 & 0.00356458 & 0.00239408 & 0.0227156 & 0.0149031 & 0.265709 & 0.262351 & -0.0126384 & 9 & 10 \\
esdown & 50000 & all & $\tau_4^{\mathrm{WTA}}$ & 50000 & -0.0233291 & -0.019651 & 0.0497941 & 0.0365929 & 0.910554 & 0.907008 & -0.00389401 & 1 & 1 \\
esdown & 50000 & all & $\sqrt{d_{12}}$ & 50000 & -0.0153948 & -0.0112218 & 0.0386685 & 0.0277771 & 0.263809 & 0.263608 & -0.00076243 & 10 & 9 \\
esdown & 50000 & all & $\sqrt{d_{23}}$ & 50000 & -0.014384 & -0.00982974 & 0.0441975 & 0.0306493 & 0.380904 & 0.380689 & -0.000562554 & 5 & 5 \\
esdown & 50000 & all & ECF$_1$ & 50000 & 0.00954055 & 0.0107554 & 0.0504983 & 0.0293943 & 0.152917 & 0.156631 & 0.0242857 & 12 & 12 \\
esdown & 50000 & all & ECF$_2$ & 50000 & -0.00143357 & -0.00146727 & 0.0162274 & 0.00983865 & 0.0443977 & 0.0446487 & 0.00565385 & 15 & 15 \\
esdown & 50000 & all & ECF$_3$ & 50000 & -0.0333813 & -0.0264992 & 0.132161 & 0.0812522 & 0.7065 & 0.711993 & 0.00777552 & 2 & 2 \\
esdown & 50000 & all & $C_2$ & 50000 & -0.000136498 & 7.31051e-05 & 0.0107791 & 0.00668604 & 0.0957288 & 0.0960016 & 0.00284905 & 14 & 14 \\
esdown & 50000 & all & $D_2$ & 50000 & 0.0120081 & 0.00947138 & 0.0247866 & 0.0173726 & 0.288376 & 0.293833 & 0.0189246 & 7 & 6 \\
esdown & 50000 & all & $L_2$ & 50000 & -0.00422468 & -0.000689641 & 0.0480584 & 0.0304961 & 0.41236 & 0.390877 & -0.052099 & 3 & 4 \\
esdown & 50000 & all & $L_3$ & 50000 & -0.00525858 & -0.00483424 & 0.0222177 & 0.0146367 & 0.138349 & 0.142915 & 0.0330005 & 13 & 13 \\
esdown & 50000 & all & $Q_w$ & 50000 & -0.0261662 & -0.0168796 & 0.0442678 & 0.0287176 & 0.272777 & 0.26623 & -0.0240011 & 8 & 8 \\
esdown & 50000 & all & ThrustMaj & 50000 & 0.00464036 & 0.00275427 & 0.046718 & 0.032909 & 0.41146 & 0.403564 & -0.0191904 & 4 & 3 \\
esdown & 50000 & signal & $\tau_1^{\mathrm{WTA}}$ & 25927 & 0.00769006 & 0.00659946 & 0.0185624 & 0.0134839 & 0.199567 & 0.207069 & 0.0375914 & 10 & 10 \\
esdown & 50000 & signal & $\tau_2^{\mathrm{WTA}}$ & 25927 & 0.00816798 & 0.00619876 & 0.0198572 & 0.0148861 & 0.260428 & 0.257596 & -0.0108707 & 8 & 8 \\
esdown & 50000 & signal & $\tau_3^{\mathrm{WTA}}$ & 25927 & 0.00345114 & 0.00279729 & 0.0249613 & 0.0169246 & 0.295006 & 0.291071 & -0.0133372 & 6 & 6 \\
esdown & 50000 & signal & $\tau_4^{\mathrm{WTA}}$ & 25927 & -0.0288974 & -0.0246541 & 0.0515973 & 0.0389872 & 0.874644 & 0.872663 & -0.00226467 & 1 & 1 \\
esdown & 50000 & signal & $\sqrt{d_{12}}$ & 25927 & -0.0174792 & -0.0141068 & 0.0386125 & 0.0284778 & 0.196997 & 0.187092 & -0.0502847 & 11 & 11 \\
esdown & 50000 & signal & $\sqrt{d_{23}}$ & 25927 & -0.0148189 & -0.00926033 & 0.0480251 & 0.034622 & 0.350083 & 0.343645 & -0.0183884 & 5 & 5 \\
esdown & 50000 & signal & ECF$_1$ & 25927 & 0.0166546 & 0.0134882 & 0.0491605 & 0.0330913 & 0.144833 & 0.154325 & 0.0655379 & 12 & 12 \\
esdown & 50000 & signal & ECF$_2$ & 25927 & -0.00079701 & -0.000468791 & 0.0135842 & 0.00910597 & 0.0344003 & 0.0336873 & -0.0207253 & 15 & 15 \\
esdown & 50000 & signal & ECF$_3$ & 25927 & -0.0335321 & -0.0153356 & 0.167339 & 0.110229 & 0.625481 & 0.609972 & -0.0247951 & 2 & 2 \\
esdown & 50000 & signal & $C_2$ & 25927 & -0.000373882 & -0.00015305 & 0.00974878 & 0.00639381 & 0.0803986 & 0.0809436 & 0.00677876 & 14 & 14 \\
esdown & 50000 & signal & $D_2$ & 25927 & 0.0115185 & 0.0106737 & 0.0233976 & 0.0176618 & 0.207771 & 0.220534 & 0.0614291 & 9 & 9 \\
esdown & 50000 & signal & $L_2$ & 25927 & -0.0168596 & -0.00748497 & 0.0580064 & 0.0382079 & 0.404956 & 0.375733 & -0.0721633 & 4 & 4 \\
esdown & 50000 & signal & $L_3$ & 25927 & -0.00656882 & -0.00604007 & 0.0186978 & 0.0128753 & 0.0827248 & 0.0898128 & 0.0856819 & 13 & 13 \\
esdown & 50000 & signal & $Q_w$ & 25927 & -0.0344572 & -0.0259801 & 0.0560514 & 0.0407192 & 0.289637 & 0.269433 & -0.069756 & 7 & 7 \\
esdown & 50000 & signal & ThrustMaj & 25927 & 0.00342555 & 0.00117803 & 0.054437 & 0.0414562 & 0.446439 & 0.438389 & -0.0180325 & 3 & 3 \\
esdown & 50000 & background & $\tau_1^{\mathrm{WTA}}$ & 24073 & 0.00303566 & 0.00314736 & 0.0212826 & 0.0137193 & 0.28309 & 0.285506 & 0.00853656 & 9 & 9 \\
esdown & 50000 & background & $\tau_2^{\mathrm{WTA}}$ & 24073 & 0.00637364 & 0.00406823 & 0.0206776 & 0.0139463 & 0.319964 & 0.315288 & -0.0146141 & 8 & 8 \\
esdown & 50000 & background & $\tau_3^{\mathrm{WTA}}$ & 24073 & 0.00368674 & 0.00211626 & 0.020297 & 0.0128492 & 0.234154 & 0.231417 & -0.0116899 & 11 & 11 \\
esdown & 50000 & background & $\tau_4^{\mathrm{WTA}}$ & 24073 & -0.0173321 & -0.0149391 & 0.0478521 & 0.0339202 & 0.949233 & 0.94399 & -0.00552297 & 1 & 1 \\
esdown & 50000 & background & $\sqrt{d_{12}}$ & 24073 & -0.0131498 & -0.00803709 & 0.0387288 & 0.0269999 & 0.335768 & 0.346017 & 0.0305249 & 7 & 7 \\
esdown & 50000 & background & $\sqrt{d_{23}}$ & 24073 & -0.0139156 & -0.010448 & 0.0400752 & 0.0270003 & 0.414102 & 0.420583 & 0.0156508 & 4 & 3 \\
esdown & 50000 & background & ECF$_1$ & 24073 & 0.0018786 & 0.00825118 & 0.0519391 & 0.0258905 & 0.161622 & 0.159115 & -0.0155166 & 13 & 13 \\
esdown & 50000 & background & ECF$_2$ & 24073 & -0.00211916 & -0.0027123 & 0.0190743 & 0.010774 & 0.0551649 & 0.0564539 & 0.023367 & 15 & 15 \\
esdown & 50000 & background & ECF$_3$ & 24073 & -0.0332189 & -0.0352516 & 0.0942744 & 0.0648063 & 0.793759 & 0.821867 & 0.0354114 & 2 & 2 \\
esdown & 50000 & background & $C_2$ & 24073 & 0.000119169 & 0.000328286 & 0.0118889 & 0.00705153 & 0.112241 & 0.11222 & -0.000189183 & 14 & 14 \\
esdown & 50000 & background & $D_2$ & 24073 & 0.0125355 & 0.00801361 & 0.0262825 & 0.0170362 & 0.375188 & 0.372779 & -0.00642113 & 5 & 5 \\
esdown & 50000 & background & $L_2$ & 24073 & 0.00938329 & 0.00381768 & 0.0373443 & 0.0241438 & 0.420332 & 0.407186 & -0.0312762 & 3 & 4 \\
esdown & 50000 & background & $L_3$ & 24073 & -0.00384742 & -0.00311358 & 0.0260088 & 0.0171657 & 0.198258 & 0.200104 & 0.00930996 & 12 & 12 \\
esdown & 50000 & background & $Q_w$ & 24073 & -0.0172367 & -0.0116993 & 0.0315767 & 0.0204435 & 0.254618 & 0.262781 & 0.0320604 & 10 & 10 \\
esdown & 50000 & background & ThrustMaj & 24073 & 0.00594873 & 0.00374602 & 0.0384045 & 0.0254426 & 0.373784 & 0.366059 & -0.0206671 & 6 & 6 \\
esdown & 100000 & all & $\tau_1^{\mathrm{WTA}}$ & 100000 & 0.00534435 & 0.00486714 & 0.0198508 & 0.0135393 & 0.239359 & 0.244408 & 0.0210942 & 11 & 11 \\
esdown & 100000 & all & $\tau_2^{\mathrm{WTA}}$ & 100000 & 0.00720541 & 0.00506014 & 0.0202738 & 0.014444 & 0.289806 & 0.285956 & -0.0132858 & 6 & 7 \\
esdown & 100000 & all & $\tau_3^{\mathrm{WTA}}$ & 100000 & 0.00346677 & 0.00234497 & 0.0228598 & 0.0149684 & 0.266534 & 0.263137 & -0.0127462 & 9 & 9 \\
esdown & 100000 & all & $\tau_4^{\mathrm{WTA}}$ & 100000 & -0.0232354 & -0.0193963 & 0.0498115 & 0.0366069 & 0.910574 & 0.907307 & -0.00358751 & 1 & 1 \\
esdown & 100000 & all & $\sqrt{d_{12}}$ & 100000 & -0.0156328 & -0.0113007 & 0.0389934 & 0.0279073 & 0.262793 & 0.262643 & -0.000568505 & 10 & 10 \\
esdown & 100000 & all & $\sqrt{d_{23}}$ & 100000 & -0.0141728 & -0.00979245 & 0.0442657 & 0.0306614 & 0.381419 & 0.380805 & -0.00160959 & 5 & 5 \\
esdown & 100000 & all & ECF$_1$ & 100000 & 0.00971306 & 0.0109327 & 0.0504284 & 0.0293521 & 0.15328 & 0.156883 & 0.0235102 & 12 & 12 \\
esdown & 100000 & all & ECF$_2$ & 100000 & -0.00146359 & -0.00142894 & 0.0160902 & 0.00973163 & 0.0442419 & 0.0443198 & 0.00175933 & 15 & 15 \\
esdown & 100000 & all & ECF$_3$ & 100000 & -0.0328923 & -0.0262338 & 0.132332 & 0.0812167 & 0.707233 & 0.71297 & 0.00811114 & 2 & 2 \\
esdown & 100000 & all & $C_2$ & 100000 & -8.99549e-05 & 6.1186e-05 & 0.0108126 & 0.00672704 & 0.0957582 & 0.0959462 & 0.00196282 & 14 & 14 \\
esdown & 100000 & all & $D_2$ & 100000 & 0.0119329 & 0.00941069 & 0.0247884 & 0.0173063 & 0.286963 & 0.292475 & 0.0192094 & 7 & 6 \\
esdown & 100000 & all & $L_2$ & 100000 & -0.00401339 & -0.000528201 & 0.0480699 & 0.0306167 & 0.412814 & 0.391654 & -0.0512574 & 3 & 4 \\
esdown & 100000 & all & $L_3$ & 100000 & -0.00529391 & -0.00491074 & 0.0221455 & 0.0145749 & 0.138091 & 0.142682 & 0.0332463 & 13 & 13 \\
esdown & 100000 & all & $Q_w$ & 100000 & -0.0265356 & -0.017056 & 0.0446582 & 0.0288031 & 0.272023 & 0.265313 & -0.0246656 & 8 & 8 \\
esdown & 100000 & all & ThrustMaj & 100000 & 0.00433202 & 0.0025569 & 0.0467127 & 0.0329132 & 0.410402 & 0.402607 & -0.0189939 & 4 & 3 \\
esdown & 100000 & signal & $\tau_1^{\mathrm{WTA}}$ & 51853 & 0.00767254 & 0.00655052 & 0.0186201 & 0.0135044 & 0.199249 & 0.20682 & 0.0379976 & 10 & 10 \\
esdown & 100000 & signal & $\tau_2^{\mathrm{WTA}}$ & 51853 & 0.00804523 & 0.00616318 & 0.0198201 & 0.0148828 & 0.260987 & 0.258065 & -0.0111965 & 8 & 8 \\
esdown & 100000 & signal & $\tau_3^{\mathrm{WTA}}$ & 51853 & 0.00337349 & 0.00272498 & 0.0250432 & 0.0170673 & 0.295619 & 0.291542 & -0.0137899 & 6 & 6 \\
esdown & 100000 & signal & $\tau_4^{\mathrm{WTA}}$ & 51853 & -0.0287042 & -0.0241066 & 0.0516104 & 0.0388809 & 0.876644 & 0.875072 & -0.00179369 & 1 & 1 \\
esdown & 100000 & signal & $\sqrt{d_{12}}$ & 51853 & -0.0179582 & -0.0144712 & 0.0389669 & 0.02857 & 0.196815 & 0.186723 & -0.0512773 & 11 & 11 \\
esdown & 100000 & signal & $\sqrt{d_{23}}$ & 51853 & -0.0145739 & -0.00937439 & 0.048091 & 0.0346517 & 0.350571 & 0.343429 & -0.0203709 & 5 & 5 \\
esdown & 100000 & signal & ECF$_1$ & 51853 & 0.0170519 & 0.0136396 & 0.0489855 & 0.0330192 & 0.145377 & 0.155069 & 0.0666689 & 12 & 12 \\
esdown & 100000 & signal & ECF$_2$ & 51853 & -0.000838902 & -0.000469965 & 0.0134496 & 0.00895782 & 0.0343047 & 0.0334075 & -0.0261551 & 15 & 15 \\
esdown & 100000 & signal & ECF$_3$ & 51853 & -0.0331425 & -0.0153742 & 0.167648 & 0.109858 & 0.62518 & 0.610168 & -0.0240124 & 2 & 2 \\
esdown & 100000 & signal & $C_2$ & 51853 & -0.000398099 & -0.000163279 & 0.00977628 & 0.00641411 & 0.079899 & 0.080366 & 0.00584547 & 14 & 14 \\
esdown & 100000 & signal & $D_2$ & 51853 & 0.0114525 & 0.0107328 & 0.0234113 & 0.0176016 & 0.207521 & 0.220392 & 0.0620253 & 9 & 9 \\
esdown & 100000 & signal & $L_2$ & 51853 & -0.016591 & -0.00715986 & 0.0579745 & 0.0384137 & 0.406232 & 0.377578 & -0.0705357 & 4 & 4 \\
esdown & 100000 & signal & $L_3$ & 51853 & -0.00659795 & -0.00614977 & 0.018527 & 0.0127346 & 0.0822517 & 0.0891937 & 0.0843994 & 13 & 13 \\
esdown & 100000 & signal & $Q_w$ & 51853 & -0.0347862 & -0.0260493 & 0.056534 & 0.0410207 & 0.288838 & 0.268278 & -0.0711832 & 7 & 7 \\
esdown & 100000 & signal & ThrustMaj & 51853 & 0.00295539 & 0.0008533 & 0.0542549 & 0.0414427 & 0.445973 & 0.43795 & -0.0179895 & 3 & 3 \\
esdown & 100000 & background & $\tau_1^{\mathrm{WTA}}$ & 48147 & 0.00283696 & 0.00301856 & 0.0211763 & 0.0135757 & 0.282555 & 0.28489 & 0.00826286 & 9 & 9 \\
esdown & 100000 & background & $\tau_2^{\mathrm{WTA}}$ & 48147 & 0.00630096 & 0.00398201 & 0.0207624 & 0.013987 & 0.320846 & 0.315998 & -0.0151097 & 8 & 8 \\
esdown & 100000 & background & $\tau_3^{\mathrm{WTA}}$ & 48147 & 0.00356723 & 0.00207207 & 0.0205082 & 0.0128736 & 0.235213 & 0.232549 & -0.0113271 & 11 & 11 \\
esdown & 100000 & background & $\tau_4^{\mathrm{WTA}}$ & 48147 & -0.0173456 & -0.0148017 & 0.0478742 & 0.0341155 & 0.94713 & 0.942027 & -0.00538772 & 1 & 1 \\
esdown & 100000 & background & $\sqrt{d_{12}}$ & 48147 & -0.0131284 & -0.00791937 & 0.0390219 & 0.027162 & 0.333853 & 0.344414 & 0.0316328 & 7 & 7 \\
esdown & 100000 & background & $\sqrt{d_{23}}$ & 48147 & -0.0137408 & -0.010205 & 0.0401459 & 0.0268919 & 0.414641 & 0.421054 & 0.0154653 & 4 & 3 \\
esdown & 100000 & background & ECF$_1$ & 48147 & 0.00180934 & 0.00845207 & 0.0519824 & 0.0257647 & 0.16179 & 0.158837 & -0.0182502 & 13 & 13 \\
esdown & 100000 & background & ECF$_2$ & 48147 & -0.00213637 & -0.00263203 & 0.018934 & 0.0107668 & 0.0549448 & 0.0560721 & 0.0205169 & 15 & 15 \\
esdown & 100000 & background & ECF$_3$ & 48147 & -0.0326229 & -0.0343883 & 0.0942976 & 0.06462 & 0.795605 & 0.823693 & 0.0353034 & 2 & 2 \\
esdown & 100000 & background & $C_2$ & 48147 & 0.000241908 & 0.000307426 & 0.0119288 & 0.00712524 & 0.112839 & 0.112726 & -0.00100495 & 14 & 14 \\
esdown & 100000 & background & $D_2$ & 48147 & 0.0124503 & 0.0078131 & 0.0262715 & 0.0169691 & 0.372517 & 0.370102 & -0.00648444 & 5 & 5 \\
esdown & 100000 & background & $L_2$ & 48147 & 0.00953237 & 0.00381666 & 0.0374029 & 0.0241723 & 0.419905 & 0.406817 & -0.0311691 & 3 & 4 \\
esdown & 100000 & background & $L_3$ & 48147 & -0.00388949 & -0.00313222 & 0.0260424 & 0.0172283 & 0.198228 & 0.200289 & 0.0103945 & 12 & 12 \\
esdown & 100000 & background & $Q_w$ & 48147 & -0.0176499 & -0.0118033 & 0.0318683 & 0.0203964 & 0.253911 & 0.262117 & 0.0323205 & 10 & 10 \\
esdown & 100000 & background & ThrustMaj & 48147 & 0.00581462 & 0.0035831 & 0.03859 & 0.0254574 & 0.372097 & 0.364547 & -0.0202899 & 6 & 6 \\
cpos & 25000 & all & $\tau_1^{\mathrm{WTA}}$ & 25000 & 3.42933e-05 & 0 & 0.00294052 & 3.18326e-05 & 0.239263 & 0.239344 & 0.000337181 & 11 & 11 \\
cpos & 25000 & all & $\tau_2^{\mathrm{WTA}}$ & 25000 & 0.000205814 & 0 & 0.00453562 & 3.68655e-05 & 0.290105 & 0.290034 & -0.000244393 & 6 & 6 \\
cpos & 25000 & all & $\tau_3^{\mathrm{WTA}}$ & 25000 & -4.36753e-05 & 0 & 0.00555819 & 6.50883e-05 & 0.266625 & 0.266693 & 0.000253173 & 9 & 9 \\
cpos & 25000 & all & $\tau_4^{\mathrm{WTA}}$ & 25000 & -0.000414701 & 0 & 0.0131249 & 0.000388533 & 0.907388 & 0.907078 & -0.00034217 & 1 & 1 \\
cpos & 25000 & all & $\sqrt{d_{12}}$ & 25000 & -9.57921e-06 & 0 & 0.00543589 & 0.000150882 & 0.264572 & 0.264453 & -0.000450462 & 10 & 10 \\
cpos & 25000 & all & $\sqrt{d_{23}}$ & 25000 & -3.68057e-05 & 0 & 0.00801595 & 0.000283554 & 0.381048 & 0.380997 & -0.000134602 & 5 & 5 \\
cpos & 25000 & all & ECF$_1$ & 25000 & -1.08697e-05 & 0 & 0.00319144 & 7.10189e-05 & 0.15373 & 0.153702 & -0.000181551 & 12 & 12 \\
cpos & 25000 & all & ECF$_2$ & 25000 & -4.37941e-05 & 0 & 0.00166235 & 1.27628e-05 & 0.0442467 & 0.044218 & -0.000648039 & 15 & 15 \\
cpos & 25000 & all & ECF$_3$ & 25000 & 5.28998e-05 & 0 & 0.00843733 & 0.000395387 & 0.715047 & 0.715101 & 7.53553e-05 & 2 & 2 \\
cpos & 25000 & all & $C_2$ & 25000 & 2.08995e-05 & 0 & 0.00214031 & 3.7428e-05 & 0.0974392 & 0.0974762 & 0.000379108 & 14 & 14 \\
cpos & 25000 & all & $D_2$ & 25000 & -0.000126296 & 0 & 0.00402753 & 0.000115213 & 0.285922 & 0.285959 & 0.00013029 & 7 & 7 \\
cpos & 25000 & all & $L_2$ & 25000 & 0.0003422 & 0 & 0.00624741 & 0.000269104 & 0.415616 & 0.415412 & -0.000488751 & 3 & 3 \\
cpos & 25000 & all & $L_3$ & 25000 & -6.88588e-05 & 0 & 0.00348623 & 6.9987e-05 & 0.139558 & 0.139489 & -0.00049906 & 13 & 13 \\
cpos & 25000 & all & $Q_w$ & 25000 & 0.000337506 & 0 & 0.00603792 & 0.000174262 & 0.27107 & 0.271326 & 0.000943313 & 8 & 8 \\
cpos & 25000 & all & ThrustMaj & 25000 & 0.000369438 & 0 & 0.00737067 & 0.000285715 & 0.411984 & 0.411954 & -7.23385e-05 & 4 & 4 \\
cpos & 25000 & signal & $\tau_1^{\mathrm{WTA}}$ & 12963 & -4.65354e-05 & 0 & 0.00230213 & 3.57628e-07 & 0.199802 & 0.199744 & -0.000292054 & 10 & 10 \\
cpos & 25000 & signal & $\tau_2^{\mathrm{WTA}}$ & 12963 & 0.000193769 & 0 & 0.00348952 & 1.53482e-06 & 0.261553 & 0.261412 & -0.000540435 & 8 & 8 \\
cpos & 25000 & signal & $\tau_3^{\mathrm{WTA}}$ & 12963 & -7.30481e-05 & 0 & 0.00504036 & 9.0003e-06 & 0.296385 & 0.296514 & 0.000437204 & 6 & 6 \\
cpos & 25000 & signal & $\tau_4^{\mathrm{WTA}}$ & 12963 & -0.000823612 & 0 & 0.01111 & 8.51452e-05 & 0.876649 & 0.876492 & -0.000178954 & 1 & 1 \\
cpos & 25000 & signal & $\sqrt{d_{12}}$ & 12963 & 4.19361e-05 & 0 & 0.00405804 & 3.2872e-05 & 0.195527 & 0.195363 & -0.000841285 & 11 & 11 \\
cpos & 25000 & signal & $\sqrt{d_{23}}$ & 12963 & 0.000158237 & 0 & 0.00615944 & 7.48932e-05 & 0.349928 & 0.349846 & -0.000232591 & 5 & 5 \\
cpos & 25000 & signal & ECF$_1$ & 12963 & -4.31137e-05 & 0 & 0.00221323 & 1.13249e-05 & 0.144886 & 0.144748 & -0.000953395 & 12 & 12 \\
cpos & 25000 & signal & ECF$_2$ & 12963 & -8.52821e-06 & 0 & 0.00108731 & 1.49012e-08 & 0.0339882 & 0.0339349 & -0.00156955 & 15 & 15 \\
cpos & 25000 & signal & ECF$_3$ & 12963 & 7.99052e-05 & 0 & 0.00718308 & 0.00011456 & 0.625401 & 0.625362 & -6.36645e-05 & 2 & 2 \\
cpos & 25000 & signal & $C_2$ & 12963 & -3.0085e-05 & 0 & 0.00164417 & 3.76254e-06 & 0.0800476 & 0.0801223 & 0.000932723 & 14 & 14 \\
cpos & 25000 & signal & $D_2$ & 12963 & -9.14908e-05 & 0 & 0.00274349 & 2.23368e-05 & 0.207138 & 0.20712 & -8.70453e-05 & 9 & 9 \\
cpos & 25000 & signal & $L_2$ & 12963 & 8.89963e-05 & 0 & 0.00532383 & 7.04527e-05 & 0.406923 & 0.40679 & -0.000328693 & 4 & 4 \\
cpos & 25000 & signal & $L_3$ & 12963 & -0.000113422 & 0 & 0.00219947 & 1.20979e-05 & 0.081696 & 0.0816449 & -0.000626353 & 13 & 13 \\
cpos & 25000 & signal & $Q_w$ & 12963 & 0.000381205 & 0 & 0.0052849 & 4.46439e-05 & 0.28755 & 0.287619 & 0.000239725 & 7 & 7 \\
cpos & 25000 & signal & ThrustMaj & 12963 & 0.000248874 & 0 & 0.00651146 & 7.07507e-05 & 0.449377 & 0.44955 & 0.000384519 & 3 & 3 \\
cpos & 25000 & background & $\tau_1^{\mathrm{WTA}}$ & 12037 & 0.00012134 & 0 & 0.00362801 & 0.000142813 & 0.281762 & 0.281989 & 0.000806823 & 9 & 9 \\
cpos & 25000 & background & $\tau_2^{\mathrm{WTA}}$ & 12037 & 0.000218785 & 0 & 0.0056622 & 0.000149101 & 0.320853 & 0.320857 & 1.21679e-05 & 8 & 8 \\
cpos & 25000 & background & $\tau_3^{\mathrm{WTA}}$ & 12037 & -1.20428e-05 & 0 & 0.00611585 & 0.00019446 & 0.234576 & 0.234576 & 7.62286e-07 & 11 & 11 \\
cpos & 25000 & background & $\tau_4^{\mathrm{WTA}}$ & 12037 & 2.56668e-05 & 0 & 0.0152949 & 0.00100687 & 0.940496 & 0.940019 & -0.000506372 & 1 & 1 \\
cpos & 25000 & background & $\sqrt{d_{12}}$ & 12037 & -6.50576e-05 & 0 & 0.00691974 & 0.000408798 & 0.338931 & 0.338857 & -0.000216045 & 7 & 7 \\
cpos & 25000 & background & $\sqrt{d_{23}}$ & 12037 & -0.000246853 & 0 & 0.0100153 & 0.000774622 & 0.414559 & 0.414541 & -4.48589e-05 & 4 & 4 \\
cpos & 25000 & background & ECF$_1$ & 12037 & 2.38547e-05 & 0 & 0.00424491 & 0.000198483 & 0.163254 & 0.163345 & 0.000557878 & 13 & 13 \\
cpos & 25000 & background & ECF$_2$ & 12037 & -8.1773e-05 & 0 & 0.00228162 & 4.88572e-05 & 0.0552946 & 0.0552923 & -4.12989e-05 & 15 & 15 \\
cpos & 25000 & background & ECF$_3$ & 12037 & 2.38169e-05 & 0 & 0.00978807 & 0.00093329 & 0.811587 & 0.811741 & 0.000189554 & 2 & 2 \\
cpos & 25000 & background & $C_2$ & 12037 & 7.58061e-05 & 0 & 0.00267461 & 0.000130709 & 0.116169 & 0.116165 & -3.33507e-05 & 14 & 14 \\
cpos & 25000 & background & $D_2$ & 12037 & -0.000163779 & 0 & 0.00541035 & 0.000307143 & 0.370768 & 0.370864 & 0.000260029 & 6 & 6 \\
cpos & 25000 & background & $L_2$ & 12037 & 0.000614883 & 0 & 0.00724205 & 0.000640824 & 0.424974 & 0.424697 & -0.000652815 & 3 & 3 \\
cpos & 25000 & background & $L_3$ & 12037 & -2.0867e-05 & 0 & 0.00487197 & 0.000210514 & 0.201873 & 0.201783 & -0.000442741 & 12 & 12 \\
cpos & 25000 & background & $Q_w$ & 12037 & 0.000290445 & 0 & 0.00684887 & 0.000400126 & 0.253325 & 0.25378 & 0.00179619 & 10 & 10 \\
cpos & 25000 & background & ThrustMaj & 12037 & 0.000499277 & 0 & 0.00829597 & 0.000673801 & 0.371716 & 0.371466 & -0.000671145 & 5 & 5 \\
cpos & 50000 & all & $\tau_1^{\mathrm{WTA}}$ & 50000 & 2.57342e-05 & 0 & 0.00293048 & 3.60459e-05 & 0.238936 & 0.238978 & 0.000174496 & 11 & 11 \\
cpos & 50000 & all & $\tau_2^{\mathrm{WTA}}$ & 50000 & 0.000204392 & 0 & 0.00451146 & 3.85791e-05 & 0.289775 & 0.289672 & -0.000358111 & 6 & 6 \\
cpos & 50000 & all & $\tau_3^{\mathrm{WTA}}$ & 50000 & 6.38116e-05 & 0 & 0.00547029 & 6.7044e-05 & 0.266565 & 0.266544 & -7.6584e-05 & 9 & 9 \\
cpos & 50000 & all & $\tau_4^{\mathrm{WTA}}$ & 50000 & -0.000456742 & 0 & 0.0129519 & 0.000405431 & 0.907668 & 0.907559 & -0.000119844 & 1 & 1 \\
cpos & 50000 & all & $\sqrt{d_{12}}$ & 50000 & 6.43722e-05 & 0 & 0.00554501 & 0.000164002 & 0.263493 & 0.263386 & -0.00040582 & 10 & 10 \\
cpos & 50000 & all & $\sqrt{d_{23}}$ & 50000 & 3.89783e-05 & 0 & 0.00811507 & 0.000299573 & 0.381165 & 0.381169 & 9.69525e-06 & 5 & 5 \\
cpos & 50000 & all & ECF$_1$ & 50000 & -3.66697e-05 & 0 & 0.00320373 & 7.22678e-05 & 0.15381 & 0.153766 & -0.000285895 & 12 & 12 \\
cpos & 50000 & all & ECF$_2$ & 50000 & -4.22326e-06 & 0 & 0.00165457 & 1.23531e-05 & 0.0443042 & 0.0442828 & -0.000483485 & 15 & 15 \\
cpos & 50000 & all & ECF$_3$ & 50000 & -6.79538e-06 & 0 & 0.00860197 & 0.000423983 & 0.716062 & 0.715985 & -0.000106963 & 2 & 2 \\
cpos & 50000 & all & $C_2$ & 50000 & -1.17013e-05 & 0 & 0.00212609 & 3.79048e-05 & 0.0964756 & 0.0965221 & 0.000482132 & 14 & 14 \\
cpos & 50000 & all & $D_2$ & 50000 & -7.72105e-05 & 0 & 0.00401788 & 0.000123908 & 0.285454 & 0.285436 & -6.03452e-05 & 7 & 7 \\
cpos & 50000 & all & $L_2$ & 50000 & 0.000327696 & 0 & 0.00636399 & 0.000286147 & 0.413679 & 0.413439 & -0.000580372 & 3 & 3 \\
cpos & 50000 & all & $L_3$ & 50000 & -5.86798e-05 & 0 & 0.0035144 & 7.19e-05 & 0.139335 & 0.139235 & -0.000718778 & 13 & 13 \\
cpos & 50000 & all & $Q_w$ & 50000 & 0.000259844 & 0 & 0.00602829 & 0.000183668 & 0.27183 & 0.271949 & 0.00043975 & 8 & 8 \\
cpos & 50000 & all & ThrustMaj & 50000 & 0.00034003 & 0 & 0.00741499 & 0.00029403 & 0.411787 & 0.411836 & 0.000116883 & 4 & 4 \\
cpos & 50000 & signal & $\tau_1^{\mathrm{WTA}}$ & 25926 & -3.64379e-05 & 0 & 0.00232547 & 1.43982e-06 & 0.199358 & 0.199304 & -0.00027342 & 10 & 10 \\
cpos & 50000 & signal & $\tau_2^{\mathrm{WTA}}$ & 25926 & 0.000197408 & 0 & 0.00339655 & 3.13018e-06 & 0.260565 & 0.260448 & -0.000450411 & 8 & 8 \\
cpos & 50000 & signal & $\tau_3^{\mathrm{WTA}}$ & 25926 & -8.69507e-06 & 0 & 0.00495289 & 9.80496e-06 & 0.295624 & 0.295591 & -0.00011049 & 6 & 6 \\
cpos & 50000 & signal & $\tau_4^{\mathrm{WTA}}$ & 25926 & -0.000854419 & 0 & 0.0108285 & 9.98452e-05 & 0.878199 & 0.878192 & -7.94096e-06 & 1 & 1 \\
cpos & 50000 & signal & $\sqrt{d_{12}}$ & 25926 & 8.44337e-05 & 0 & 0.00406253 & 3.77018e-05 & 0.195733 & 0.195636 & -0.000493933 & 11 & 11 \\
cpos & 50000 & signal & $\sqrt{d_{23}}$ & 25926 & 0.000311541 & 0 & 0.00619883 & 8.71867e-05 & 0.350488 & 0.350459 & -8.46057e-05 & 5 & 5 \\
cpos & 50000 & signal & ECF$_1$ & 25926 & -5.35907e-05 & 0 & 0.00221538 & 1.45212e-05 & 0.145477 & 0.1454 & -0.000530585 & 12 & 12 \\
cpos & 50000 & signal & ECF$_2$ & 25926 & -2.35278e-05 & 0 & 0.0010758 & 1.63913e-07 & 0.0342281 & 0.0341629 & -0.0019052 & 15 & 15 \\
cpos & 50000 & signal & ECF$_3$ & 25926 & -0.000106259 & 0 & 0.00719283 & 0.000130504 & 0.628987 & 0.628887 & -0.000158823 & 2 & 2 \\
cpos & 50000 & signal & $C_2$ & 25926 & -5.40473e-05 & 0 & 0.00162777 & 4.62495e-06 & 0.0793848 & 0.0794512 & 0.000836333 & 14 & 14 \\
cpos & 50000 & signal & $D_2$ & 25926 & -3.92414e-05 & 0 & 0.00276067 & 2.60966e-05 & 0.208039 & 0.207996 & -0.000205999 & 9 & 9 \\
cpos & 50000 & signal & $L_2$ & 25926 & 4.38556e-05 & 0 & 0.00541999 & 8.4132e-05 & 0.405921 & 0.405799 & -0.000300798 & 4 & 4 \\
cpos & 50000 & signal & $L_3$ & 25926 & -9.19081e-05 & 0 & 0.00218118 & 1.37612e-05 & 0.081663 & 0.0816044 & -0.000716748 & 13 & 13 \\
cpos & 50000 & signal & $Q_w$ & 25926 & 0.00029302 & 0 & 0.00527156 & 5.42849e-05 & 0.288338 & 0.288369 & 0.000107287 & 7 & 7 \\
cpos & 50000 & signal & ThrustMaj & 25926 & 3.59521e-05 & 0 & 0.00642838 & 7.83056e-05 & 0.449833 & 0.450021 & 0.000418381 & 3 & 3 \\
cpos & 50000 & background & $\tau_1^{\mathrm{WTA}}$ & 24074 & 9.26892e-05 & 0 & 0.00358203 & 0.000145487 & 0.28156 & 0.281705 & 0.000513887 & 9 & 9 \\
cpos & 50000 & background & $\tau_2^{\mathrm{WTA}}$ & 24074 & 0.000211914 & 0 & 0.00571215 & 0.000141956 & 0.32123 & 0.32114 & -0.000281852 & 8 & 8 \\
cpos & 50000 & background & $\tau_3^{\mathrm{WTA}}$ & 24074 & 0.000141896 & 0 & 0.00602749 & 0.000188351 & 0.235273 & 0.235264 & -3.82546e-05 & 11 & 11 \\
cpos & 50000 & background & $\tau_4^{\mathrm{WTA}}$ & 24074 & -2.8471e-05 & 0 & 0.0152386 & 0.00101373 & 0.939414 & 0.939194 & -0.000233999 & 1 & 1 \\
cpos & 50000 & background & $\sqrt{d_{12}}$ & 24074 & 4.27673e-05 & 0 & 0.00714155 & 0.000422537 & 0.336466 & 0.336349 & -0.000347567 & 7 & 7 \\
cpos & 50000 & background & $\sqrt{d_{23}}$ & 24074 & -0.000254552 & 0 & 0.0101787 & 0.000755869 & 0.4142 & 0.414244 & 0.000104689 & 4 & 4 \\
cpos & 50000 & background & ECF$_1$ & 24074 & -1.84471e-05 & 0 & 0.00426812 & 0.000198789 & 0.162783 & 0.162775 & -5.08961e-05 & 13 & 13 \\
cpos & 50000 & background & ECF$_2$ & 24074 & 1.65663e-05 & 0 & 0.00227785 & 4.75952e-05 & 0.0551557 & 0.0551816 & 0.000469682 & 15 & 15 \\
cpos & 50000 & background & ECF$_3$ & 24074 & 0.00010032 & 0 & 0.0101195 & 0.000954002 & 0.809836 & 0.809783 & -6.602e-05 & 2 & 2 \\
cpos & 50000 & background & $C_2$ & 24074 & 3.39024e-05 & 0 & 0.00266274 & 0.000125099 & 0.114881 & 0.114906 & 0.000222322 & 14 & 14 \\
cpos & 50000 & background & $D_2$ & 24074 & -0.000118101 & 0 & 0.0053718 & 0.000322171 & 0.368824 & 0.368838 & 3.64424e-05 & 6 & 6 \\
cpos & 50000 & background & $L_2$ & 24074 & 0.000633371 & 0 & 0.00738061 & 0.000664087 & 0.422037 & 0.421667 & -0.000877537 & 3 & 3 \\
cpos & 50000 & background & $L_3$ & 24074 & -2.28953e-05 & 0 & 0.00495019 & 0.000207413 & 0.201442 & 0.201297 & -0.000718126 & 12 & 12 \\
cpos & 50000 & background & $Q_w$ & 24074 & 0.000224117 & 0 & 0.00684324 & 0.000402324 & 0.254051 & 0.254267 & 0.000850018 & 10 & 10 \\
cpos & 50000 & background & ThrustMaj & 24074 & 0.000667501 & 0 & 0.00847751 & 0.000677586 & 0.370813 & 0.370712 & -0.000270607 & 5 & 5 \\
cpos & 100000 & all & $\tau_1^{\mathrm{WTA}}$ & 100000 & 2.10144e-05 & 0 & 0.00288102 & 3.61837e-05 & 0.23918 & 0.239186 & 2.42974e-05 & 11 & 11 \\
cpos & 100000 & all & $\tau_2^{\mathrm{WTA}}$ & 100000 & 0.000154604 & 0 & 0.00448118 & 3.66122e-05 & 0.289843 & 0.289711 & -0.000455811 & 6 & 6 \\
cpos & 100000 & all & $\tau_3^{\mathrm{WTA}}$ & 100000 & 6.09334e-05 & 0 & 0.0055386 & 6.38273e-05 & 0.266878 & 0.266916 & 0.00014171 & 9 & 9 \\
cpos & 100000 & all & $\tau_4^{\mathrm{WTA}}$ & 100000 & -0.00052115 & 0 & 0.0128198 & 0.000397295 & 0.909541 & 0.909397 & -0.000158589 & 1 & 1 \\
cpos & 100000 & all & $\sqrt{d_{12}}$ & 100000 & 9.16216e-06 & 0 & 0.00550773 & 0.000163957 & 0.263711 & 0.263627 & -0.000319031 & 10 & 10 \\
cpos & 100000 & all & $\sqrt{d_{23}}$ & 100000 & 0.000109435 & 0 & 0.00806241 & 0.000295162 & 0.38138 & 0.38141 & 7.71275e-05 & 5 & 5 \\
cpos & 100000 & all & ECF$_1$ & 100000 & -2.31914e-05 & 0 & 0.00318026 & 7.07805e-05 & 0.153813 & 0.153805 & -5.38643e-05 & 12 & 12 \\
cpos & 100000 & all & ECF$_2$ & 100000 & -3.38938e-06 & 0 & 0.00164582 & 1.20346e-05 & 0.0444089 & 0.0443927 & -0.000364988 & 15 & 15 \\
cpos & 100000 & all & ECF$_3$ & 100000 & -9.78237e-05 & 0 & 0.00849404 & 0.000418544 & 0.711873 & 0.711793 & -0.000112951 & 2 & 2 \\
cpos & 100000 & all & $C_2$ & 100000 & -1.25456e-06 & 0 & 0.00210311 & 3.66867e-05 & 0.0961373 & 0.0961621 & 0.000258538 & 14 & 14 \\
cpos & 100000 & all & $D_2$ & 100000 & -2.65281e-05 & 0 & 0.00396589 & 0.000119202 & 0.285994 & 0.285943 & -0.000178713 & 7 & 7 \\
cpos & 100000 & all & $L_2$ & 100000 & 0.000259015 & 0 & 0.00631768 & 0.00028877 & 0.414281 & 0.414111 & -0.000411626 & 3 & 3 \\
cpos & 100000 & all & $L_3$ & 100000 & -2.99949e-05 & 0 & 0.00345842 & 7.07097e-05 & 0.13873 & 0.13867 & -0.000433512 & 13 & 13 \\
cpos & 100000 & all & $Q_w$ & 100000 & 0.000246391 & 0 & 0.00585732 & 0.000181288 & 0.27118 & 0.271312 & 0.000485752 & 8 & 8 \\
cpos & 100000 & all & ThrustMaj & 100000 & 0.000216746 & 0 & 0.00730477 & 0.000288695 & 0.411233 & 0.411196 & -9.14579e-05 & 4 & 4 \\
cpos & 100000 & signal & $\tau_1^{\mathrm{WTA}}$ & 51852 & -3.33574e-05 & 0 & 0.00226642 & 1.04308e-06 & 0.199154 & 0.199116 & -0.00018945 & 10 & 10 \\
cpos & 100000 & signal & $\tau_2^{\mathrm{WTA}}$ & 51852 & 0.000182688 & 0 & 0.00337933 & 2.27243e-06 & 0.26106 & 0.260929 & -0.000503555 & 8 & 8 \\
cpos & 100000 & signal & $\tau_3^{\mathrm{WTA}}$ & 51852 & -3.025e-05 & 0 & 0.00509897 & 9.37283e-06 & 0.296904 & 0.296885 & -6.25349e-05 & 6 & 6 \\
cpos & 100000 & signal & $\tau_4^{\mathrm{WTA}}$ & 51852 & -0.00081385 & 0 & 0.0107635 & 9.68166e-05 & 0.877911 & 0.877768 & -0.000162741 & 1 & 1 \\
cpos & 100000 & signal & $\sqrt{d_{12}}$ & 51852 & 1.67801e-05 & 0 & 0.00397092 & 3.57628e-05 & 0.196782 & 0.19672 & -0.00031312 & 11 & 11 \\
cpos & 100000 & signal & $\sqrt{d_{23}}$ & 51852 & 0.0002832 & 0 & 0.00606789 & 7.96318e-05 & 0.350749 & 0.350844 & 0.000268837 & 5 & 5 \\
cpos & 100000 & signal & ECF$_1$ & 51852 & -2.81989e-05 & 0 & 0.00216523 & 1.21612e-05 & 0.145068 & 0.145006 & -0.000432033 & 12 & 12 \\
cpos & 100000 & signal & ECF$_2$ & 51852 & -2.61724e-05 & 0 & 0.0010738 & 6.33299e-08 & 0.0342996 & 0.0342522 & -0.00138098 & 15 & 15 \\
cpos & 100000 & signal & ECF$_3$ & 51852 & -0.000252365 & 0 & 0.00703136 & 0.000121385 & 0.626462 & 0.626303 & -0.000253942 & 2 & 2 \\
cpos & 100000 & signal & $C_2$ & 51852 & -3.68573e-05 & 0 & 0.00159958 & 3.76254e-06 & 0.0798236 & 0.0798782 & 0.000683048 & 14 & 14 \\
cpos & 100000 & signal & $D_2$ & 51852 & -5.31443e-05 & 0 & 0.00271412 & 2.48253e-05 & 0.207442 & 0.2074 & -0.000203503 & 9 & 9 \\
cpos & 100000 & signal & $L_2$ & 51852 & 0.000179278 & 0 & 0.00540132 & 7.59363e-05 & 0.406223 & 0.406185 & -9.35398e-05 & 4 & 4 \\
cpos & 100000 & signal & $L_3$ & 51852 & -4.4896e-05 & 0 & 0.00214467 & 1.22823e-05 & 0.0817971 & 0.0817803 & -0.000205581 & 13 & 13 \\
cpos & 100000 & signal & $Q_w$ & 51852 & 0.000264986 & 0 & 0.00510033 & 5.00083e-05 & 0.287939 & 0.288131 & 0.000665415 & 7 & 7 \\
cpos & 100000 & signal & ThrustMaj & 51852 & 6.73802e-06 & 0 & 0.00630085 & 7.34627e-05 & 0.448183 & 0.448263 & 0.000178342 & 3 & 3 \\
cpos & 100000 & background & $\tau_1^{\mathrm{WTA}}$ & 48148 & 7.95689e-05 & 0 & 0.00354291 & 0.000153966 & 0.282287 & 0.282339 & 0.000182116 & 9 & 9 \\
cpos & 100000 & background & $\tau_2^{\mathrm{WTA}}$ & 48148 & 0.00012436 & 0 & 0.0056678 & 0.000141561 & 0.320838 & 0.320711 & -0.000394593 & 8 & 8 \\
cpos & 100000 & background & $\tau_3^{\mathrm{WTA}}$ & 48148 & 0.000159132 & 0 & 0.00601205 & 0.000182901 & 0.23454 & 0.23464 & 0.000427645 & 11 & 11 \\
cpos & 100000 & background & $\tau_4^{\mathrm{WTA}}$ & 48148 & -0.000205932 & 0 & 0.0150342 & 0.00101137 & 0.943602 & 0.943465 & -0.000144905 & 1 & 1 \\
cpos & 100000 & background & $\sqrt{d_{12}}$ & 48148 & 9.58215e-07 & 0 & 0.00716276 & 0.000432909 & 0.335795 & 0.335685 & -0.000327405 & 7 & 7 \\
cpos & 100000 & background & $\sqrt{d_{23}}$ & 48148 & -7.76986e-05 & 0 & 0.0102104 & 0.000743613 & 0.414376 & 0.414331 & -0.000107162 & 4 & 4 \\
cpos & 100000 & background & ECF$_1$ & 48148 & -1.77987e-05 & 0 & 0.00427338 & 0.000207875 & 0.163231 & 0.163282 & 0.000310382 & 13 & 13 \\
cpos & 100000 & background & ECF$_2$ & 48148 & 2.11463e-05 & 0 & 0.00226184 & 4.98742e-05 & 0.0552959 & 0.0553128 & 0.000305186 & 15 & 15 \\
cpos & 100000 & background & ECF$_3$ & 48148 & 6.86069e-05 & 0 & 0.0100692 & 0.000971451 & 0.803849 & 0.803852 & 4.67139e-06 & 2 & 2 \\
cpos & 100000 & background & $C_2$ & 48148 & 3.70871e-05 & 0 & 0.00264538 & 0.000130637 & 0.113707 & 0.113698 & -7.36496e-05 & 14 & 14 \\
cpos & 100000 & background & $D_2$ & 48148 & 2.13572e-06 & 0 & 0.00531395 & 0.000322342 & 0.37059 & 0.370529 & -0.000163089 & 6 & 6 \\
cpos & 100000 & background & $L_2$ & 48148 & 0.000344886 & 0 & 0.00730454 & 0.000675935 & 0.422957 & 0.422652 & -0.000720261 & 3 & 3 \\
cpos & 100000 & background & $L_3$ & 48148 & -1.39474e-05 & 0 & 0.00487324 & 0.000220345 & 0.200043 & 0.199937 & -0.000529473 & 12 & 12 \\
cpos & 100000 & background & $Q_w$ & 48148 & 0.000226365 & 0 & 0.00667254 & 0.000401363 & 0.253131 & 0.253199 & 0.000270908 & 10 & 10 \\
cpos & 100000 & background & ThrustMaj & 48148 & 0.00044291 & 0 & 0.00838591 & 0.000678167 & 0.371439 & 0.371277 & -0.00043768 & 5 & 5 \\
cer & 25000 & all & $\tau_1^{\mathrm{WTA}}$ & 25000 & 0.000181344 & 1.86265e-06 & 0.0195862 & 0.0121055 & 0.238806 & 0.238636 & -0.000710034 & 11 & 11 \\
cer & 25000 & all & $\tau_2^{\mathrm{WTA}}$ & 25000 & 5.8264e-05 & 6.15418e-06 & 0.0238799 & 0.0161765 & 0.290322 & 0.290047 & -0.000945124 & 6 & 6 \\
cer & 25000 & all & $\tau_3^{\mathrm{WTA}}$ & 25000 & -0.00018848 & -5.44321e-05 & 0.025265 & 0.0151189 & 0.267398 & 0.267735 & 0.00126132 & 9 & 9 \\
cer & 25000 & all & $\tau_4^{\mathrm{WTA}}$ & 25000 & 0.00119553 & 0.001414 & 0.0586429 & 0.0409945 & 0.917446 & 0.917833 & 0.000422617 & 1 & 1 \\
cer & 25000 & all & $\sqrt{d_{12}}$ & 25000 & -0.00039206 & -0.000411381 & 0.0316468 & 0.0194356 & 0.264698 & 0.264712 & 5.3255e-05 & 10 & 10 \\
cer & 25000 & all & $\sqrt{d_{23}}$ & 25000 & -0.00105627 & -0.000336617 & 0.0382733 & 0.0231987 & 0.381434 & 0.381353 & -0.000214395 & 5 & 5 \\
cer & 25000 & all & ECF$_1$ & 25000 & 0.000680982 & 0.000263162 & 0.0274755 & 0.0137624 & 0.153196 & 0.153686 & 0.00320024 & 12 & 12 \\
cer & 25000 & all & ECF$_2$ & 25000 & 0.000147189 & -3.19164e-05 & 0.0109021 & 0.00617469 & 0.0441951 & 0.044228 & 0.000744045 & 15 & 15 \\
cer & 25000 & all & ECF$_3$ & 25000 & -0.00426006 & -0.000997424 & 0.061661 & 0.0372742 & 0.707695 & 0.704872 & -0.00398908 & 2 & 2 \\
cer & 25000 & all & $C_2$ & 25000 & 0.000112989 & 2.3867e-05 & 0.0107625 & 0.00652738 & 0.0969508 & 0.0968397 & -0.00114628 & 14 & 14 \\
cer & 25000 & all & $D_2$ & 25000 & 0.000194329 & 0.000213735 & 0.023789 & 0.0138367 & 0.283985 & 0.284287 & 0.00106413 & 7 & 7 \\
cer & 25000 & all & $L_2$ & 25000 & -0.00121294 & -0.000669625 & 0.0367679 & 0.0250681 & 0.414871 & 0.414113 & -0.00182706 & 3 & 3 \\
cer & 25000 & all & $L_3$ & 25000 & -6.82029e-05 & 5.15096e-05 & 0.020106 & 0.0118546 & 0.138166 & 0.138563 & 0.00287527 & 13 & 13 \\
cer & 25000 & all & $Q_w$ & 25000 & -0.000867031 & -0.000107184 & 0.0331201 & 0.0180477 & 0.272751 & 0.272253 & -0.00182441 & 8 & 8 \\
cer & 25000 & all & ThrustMaj & 25000 & -0.000904688 & -0.000527658 & 0.0394743 & 0.0265879 & 0.410391 & 0.409189 & -0.00292867 & 4 & 4 \\
cer & 25000 & signal & $\tau_1^{\mathrm{WTA}}$ & 12964 & 3.77335e-05 & 4.372e-05 & 0.0182392 & 0.0112914 & 0.198751 & 0.198692 & -0.000295398 & 11 & 10 \\
cer & 25000 & signal & $\tau_2^{\mathrm{WTA}}$ & 12964 & -0.000454372 & -0.000186215 & 0.0235228 & 0.0162596 & 0.26025 & 0.260011 & -0.000916115 & 8 & 8 \\
cer & 25000 & signal & $\tau_3^{\mathrm{WTA}}$ & 12964 & -0.000708295 & -0.000435695 & 0.0269788 & 0.0162866 & 0.298127 & 0.298724 & 0.0020041 & 6 & 6 \\
cer & 25000 & signal & $\tau_4^{\mathrm{WTA}}$ & 12964 & 0.000929266 & 0.00129583 & 0.0573109 & 0.0405031 & 0.881569 & 0.88299 & 0.00161106 & 1 & 1 \\
cer & 25000 & signal & $\sqrt{d_{12}}$ & 12964 & -0.00129329 & -0.00077017 & 0.0280172 & 0.0171252 & 0.199111 & 0.198624 & -0.00244609 & 10 & 11 \\
cer & 25000 & signal & $\sqrt{d_{23}}$ & 12964 & -0.00154506 & -0.00079985 & 0.0355023 & 0.0227038 & 0.352496 & 0.351317 & -0.00334289 & 5 & 5 \\
cer & 25000 & signal & ECF$_1$ & 12964 & -0.00033189 & 2.0586e-05 & 0.0233308 & 0.0130173 & 0.144926 & 0.145666 & 0.00510568 & 12 & 12 \\
cer & 25000 & signal & ECF$_2$ & 12964 & 3.38874e-05 & -7.32351e-05 & 0.00895104 & 0.00532294 & 0.0344277 & 0.0345036 & 0.00220351 & 15 & 15 \\
cer & 25000 & signal & ECF$_3$ & 12964 & -0.00712269 & -0.00173044 & 0.0670717 & 0.0386957 & 0.62472 & 0.619136 & -0.0089384 & 2 & 2 \\
cer & 25000 & signal & $C_2$ & 12964 & 3.38814e-05 & 6.08023e-05 & 0.00943925 & 0.00587116 & 0.0810565 & 0.0808274 & -0.00282667 & 14 & 14 \\
cer & 25000 & signal & $D_2$ & 12964 & -0.000144957 & 5.57303e-06 & 0.019452 & 0.0121417 & 0.207089 & 0.207265 & 0.000851591 & 9 & 9 \\
cer & 25000 & signal & $L_2$ & 12964 & -0.00240022 & -0.00152725 & 0.0381942 & 0.026656 & 0.408094 & 0.40652 & -0.00385844 & 4 & 4 \\
cer & 25000 & signal & $L_3$ & 12964 & -0.00033221 & -1.98837e-07 & 0.0146288 & 0.00900548 & 0.0823522 & 0.0827794 & 0.0051884 & 13 & 13 \\
cer & 25000 & signal & $Q_w$ & 12964 & -0.00179187 & -0.000418596 & 0.0353821 & 0.0200886 & 0.291381 & 0.290171 & -0.00414979 & 7 & 7 \\
cer & 25000 & signal & ThrustMaj & 12964 & -0.00204435 & -0.00105298 & 0.0420839 & 0.0289901 & 0.445813 & 0.4438 & -0.00451547 & 3 & 3 \\
cer & 25000 & background & $\tau_1^{\mathrm{WTA}}$ & 12036 & 0.000336028 & -2.46614e-05 & 0.0210371 & 0.0131029 & 0.281949 & 0.28166 & -0.00102467 & 9 & 9 \\
cer & 25000 & background & $\tau_2^{\mathrm{WTA}}$ & 12036 & 0.000610425 & 0.000158407 & 0.0242646 & 0.0161042 & 0.322716 & 0.322399 & -0.000983603 & 8 & 8 \\
cer & 25000 & background & $\tau_3^{\mathrm{WTA}}$ & 12036 & 0.000371414 & 0.000267304 & 0.023419 & 0.0139697 & 0.234301 & 0.234356 & 0.000237668 & 11 & 11 \\
cer & 25000 & background & $\tau_4^{\mathrm{WTA}}$ & 12036 & 0.00148232 & 0.00158322 & 0.0600776 & 0.0414101 & 0.956093 & 0.955361 & -0.000765808 & 1 & 1 \\
cer & 25000 & background & $\sqrt{d_{12}}$ & 12036 & 0.000578657 & 8.73841e-05 & 0.0355562 & 0.0224312 & 0.33534 & 0.335894 & 0.00165026 & 7 & 7 \\
cer & 25000 & background & $\sqrt{d_{23}}$ & 12036 & -0.000529791 & 7.0855e-05 & 0.041258 & 0.0236428 & 0.412604 & 0.413705 & 0.00266947 & 4 & 4 \\
cer & 25000 & background & ECF$_1$ & 12036 & 0.00177195 & 0.000590203 & 0.0319398 & 0.0147187 & 0.162104 & 0.162324 & 0.00135366 & 13 & 13 \\
cer & 25000 & background & ECF$_2$ & 12036 & 0.000269227 & 3.8676e-05 & 0.0130037 & 0.00726584 & 0.0547152 & 0.0547021 & -0.000238775 & 15 & 15 \\
cer & 25000 & background & ECF$_3$ & 12036 & -0.00117671 & -0.000417262 & 0.0558331 & 0.0357111 & 0.797065 & 0.797227 & 0.000203402 & 2 & 2 \\
cer & 25000 & background & $C_2$ & 12036 & 0.000198196 & -2.9739e-05 & 0.0121877 & 0.00730854 & 0.114071 & 0.114086 & 0.000134158 & 14 & 14 \\
cer & 25000 & background & $D_2$ & 12036 & 0.000559775 & 0.00052407 & 0.0284603 & 0.0164205 & 0.36681 & 0.367244 & 0.00118353 & 6 & 6 \\
cer & 25000 & background & $L_2$ & 12036 & 6.58822e-05 & -3.72976e-05 & 0.0352316 & 0.023354 & 0.422166 & 0.422292 & 0.000297129 & 3 & 3 \\
cer & 25000 & background & $L_3$ & 12036 & 0.00021616 & 0.000159446 & 0.0260056 & 0.0160744 & 0.198283 & 0.198648 & 0.00184128 & 12 & 12 \\
cer & 25000 & background & $Q_w$ & 12036 & 0.000129118 & 0.000244159 & 0.0306836 & 0.0159664 & 0.252686 & 0.252954 & 0.00106207 & 10 & 10 \\
cer & 25000 & background & ThrustMaj & 12036 & 0.000322844 & -5.25489e-05 & 0.0366636 & 0.0241381 & 0.372237 & 0.371912 & -0.000873324 & 5 & 5 \\
cer & 50000 & all & $\tau_1^{\mathrm{WTA}}$ & 50000 & 0.000108596 & 1.51917e-05 & 0.0197263 & 0.0122342 & 0.239834 & 0.239777 & -0.000238583 & 11 & 11 \\
cer & 50000 & all & $\tau_2^{\mathrm{WTA}}$ & 50000 & 0.000129714 & 3.27826e-07 & 0.0239885 & 0.0162414 & 0.290124 & 0.289725 & -0.00137484 & 6 & 6 \\
cer & 50000 & all & $\tau_3^{\mathrm{WTA}}$ & 50000 & -5.4377e-05 & 1.78814e-05 & 0.0253191 & 0.0153071 & 0.267773 & 0.267859 & 0.00032265 & 9 & 9 \\
cer & 50000 & all & $\tau_4^{\mathrm{WTA}}$ & 50000 & 0.00140543 & 0.00119638 & 0.0590088 & 0.040993 & 0.914582 & 0.914593 & 1.23826e-05 & 1 & 1 \\
cer & 50000 & all & $\sqrt{d_{12}}$ & 50000 & -0.000574548 & -0.000192262 & 0.0318575 & 0.019349 & 0.263899 & 0.263739 & -0.0006061 & 10 & 10 \\
cer & 50000 & all & $\sqrt{d_{23}}$ & 50000 & -0.000600141 & -0.00018283 & 0.0381096 & 0.023311 & 0.381286 & 0.381133 & -0.00040035 & 5 & 5 \\
cer & 50000 & all & ECF$_1$ & 50000 & 0.000839648 & 0.000359278 & 0.027422 & 0.0136793 & 0.153806 & 0.154265 & 0.00297914 & 12 & 12 \\
cer & 50000 & all & ECF$_2$ & 50000 & 0.00018847 & -1.51899e-05 & 0.011001 & 0.00617491 & 0.0442916 & 0.0443595 & 0.00153321 & 15 & 15 \\
cer & 50000 & all & ECF$_3$ & 50000 & -0.00416904 & -0.000940114 & 0.0618743 & 0.0372871 & 0.707757 & 0.70515 & -0.00368353 & 2 & 2 \\
cer & 50000 & all & $C_2$ & 50000 & 0.00018029 & 4.37079e-05 & 0.0107639 & 0.00646977 & 0.0965523 & 0.0964613 & -0.000943126 & 14 & 14 \\
cer & 50000 & all & $D_2$ & 50000 & 0.000184671 & 0.00014494 & 0.023984 & 0.013858 & 0.285472 & 0.285616 & 0.0005034 & 7 & 7 \\
cer & 50000 & all & $L_2$ & 50000 & -0.000995972 & -0.000728875 & 0.0368574 & 0.0249911 & 0.413064 & 0.412367 & -0.00168714 & 3 & 3 \\
cer & 50000 & all & $L_3$ & 50000 & -1.79546e-05 & -7.04452e-06 & 0.019963 & 0.0117931 & 0.136902 & 0.137075 & 0.00126293 & 13 & 13 \\
cer & 50000 & all & $Q_w$ & 50000 & -0.000926045 & -0.000162305 & 0.0327779 & 0.0178399 & 0.271318 & 0.270566 & -0.002771 & 8 & 8 \\
cer & 50000 & all & ThrustMaj & 50000 & -0.000765246 & -0.000474215 & 0.0393922 & 0.0264041 & 0.410588 & 0.409488 & -0.00267815 & 4 & 4 \\
cer & 50000 & signal & $\tau_1^{\mathrm{WTA}}$ & 25930 & -5.20128e-05 & 3.37586e-05 & 0.0184482 & 0.0113709 & 0.198912 & 0.199026 & 0.000569715 & 10 & 10 \\
cer & 50000 & signal & $\tau_2^{\mathrm{WTA}}$ & 25930 & -0.000393924 & -0.000167428 & 0.0235412 & 0.0163076 & 0.25919 & 0.258816 & -0.00144418 & 8 & 8 \\
cer & 50000 & signal & $\tau_3^{\mathrm{WTA}}$ & 25930 & -0.00048435 & -0.000236131 & 0.026971 & 0.0165165 & 0.297318 & 0.297453 & 0.000451768 & 6 & 6 \\
cer & 50000 & signal & $\tau_4^{\mathrm{WTA}}$ & 25930 & 0.00103681 & 0.000790656 & 0.0576229 & 0.0406968 & 0.87653 & 0.877213 & 0.000779832 & 1 & 1 \\
cer & 50000 & signal & $\sqrt{d_{12}}$ & 25930 & -0.00152333 & -0.000690246 & 0.028119 & 0.0171156 & 0.197736 & 0.197118 & -0.00312588 & 11 & 11 \\
cer & 50000 & signal & $\sqrt{d_{23}}$ & 25930 & -0.00122334 & -0.000730351 & 0.0358515 & 0.0229286 & 0.351954 & 0.350859 & -0.00311153 & 5 & 5 \\
cer & 50000 & signal & ECF$_1$ & 25930 & -9.74473e-05 & 0.000200095 & 0.0231713 & 0.0128123 & 0.14511 & 0.145816 & 0.00486694 & 12 & 12 \\
cer & 50000 & signal & ECF$_2$ & 25930 & 8.79247e-05 & -4.09938e-05 & 0.00893207 & 0.00529561 & 0.0342419 & 0.0344034 & 0.00471826 & 15 & 15 \\
cer & 50000 & signal & ECF$_3$ & 25930 & -0.00758767 & -0.00191398 & 0.0674587 & 0.0387583 & 0.625941 & 0.621249 & -0.0074967 & 2 & 2 \\
cer & 50000 & signal & $C_2$ & 25930 & 1.66754e-05 & 3.02698e-05 & 0.00944977 & 0.00586437 & 0.0799064 & 0.079832 & -0.000931761 & 14 & 14 \\
cer & 50000 & signal & $D_2$ & 25930 & -0.00012759 & 6.80238e-06 & 0.0194259 & 0.011914 & 0.206775 & 0.20701 & 0.0011374 & 9 & 9 \\
cer & 50000 & signal & $L_2$ & 25930 & -0.00229386 & -0.0014125 & 0.0381293 & 0.0263509 & 0.407999 & 0.406588 & -0.00345621 & 4 & 4 \\
cer & 50000 & signal & $L_3$ & 25930 & -0.000417158 & -8.38134e-05 & 0.0146257 & 0.00899171 & 0.0819351 & 0.0823149 & 0.00463529 & 13 & 13 \\
cer & 50000 & signal & $Q_w$ & 25930 & -0.00184646 & -0.000570099 & 0.0352425 & 0.0200367 & 0.289738 & 0.288094 & -0.00567517 & 7 & 7 \\
cer & 50000 & signal & ThrustMaj & 25930 & -0.00202524 & -0.00113289 & 0.041985 & 0.0288711 & 0.446876 & 0.444806 & -0.00463165 & 3 & 3 \\
cer & 50000 & background & $\tau_1^{\mathrm{WTA}}$ & 24070 & 0.000281616 & -2.30968e-06 & 0.0211032 & 0.0132055 & 0.283917 & 0.283679 & -0.000837542 & 9 & 9 \\
cer & 50000 & background & $\tau_2^{\mathrm{WTA}}$ & 24070 & 0.000693816 & 0.000202976 & 0.0244703 & 0.0161218 & 0.323442 & 0.32302 & -0.00130536 & 8 & 8 \\
cer & 50000 & background & $\tau_3^{\mathrm{WTA}}$ & 24070 & 0.000408821 & 0.000239246 & 0.0235396 & 0.0141202 & 0.235943 & 0.235979 & 0.000153027 & 11 & 11 \\
cer & 50000 & background & $\tau_4^{\mathrm{WTA}}$ & 24070 & 0.00180252 & 0.0015825 & 0.0605017 & 0.0413277 & 0.955577 & 0.954876 & -0.000732975 & 1 & 1 \\
cer & 50000 & background & $\sqrt{d_{12}}$ & 24070 & 0.000447555 & 0.000366077 & 0.0358848 & 0.0223795 & 0.335174 & 0.335508 & 0.00099737 & 7 & 7 \\
cer & 50000 & background & $\sqrt{d_{23}}$ & 24070 & 7.121e-05 & 0.000314899 & 0.0405422 & 0.0237489 & 0.412888 & 0.413744 & 0.00207193 & 4 & 4 \\
cer & 50000 & background & ECF$_1$ & 24070 & 0.00184916 & 0.00052565 & 0.0320011 & 0.0146886 & 0.163176 & 0.163368 & 0.0011772 & 13 & 13 \\
cer & 50000 & background & ECF$_2$ & 24070 & 0.000296784 & 1.94358e-05 & 0.0132299 & 0.0073375 & 0.0551183 & 0.0550847 & -0.00060842 & 15 & 15 \\
cer & 50000 & background & ECF$_3$ & 24070 & -0.000486243 & -7.4923e-05 & 0.0558585 & 0.0358 & 0.795892 & 0.795532 & -0.000451739 & 2 & 2 \\
cer & 50000 & background & $C_2$ & 24070 & 0.000356547 & 6.8143e-05 & 0.0121795 & 0.00724084 & 0.114485 & 0.114377 & -0.000945991 & 14 & 14 \\
cer & 50000 & background & $D_2$ & 24070 & 0.000521062 & 0.000299344 & 0.0288942 & 0.0167477 & 0.370246 & 0.370294 & 0.00013056 & 6 & 6 \\
cer & 50000 & background & $L_2$ & 24070 & 0.000402207 & -9.84436e-05 & 0.0354872 & 0.023568 & 0.418519 & 0.418592 & 0.000173679 & 3 & 3 \\
cer & 50000 & background & $L_3$ & 24070 & 0.000412098 & 0.000147335 & 0.0257128 & 0.0160707 & 0.196116 & 0.196069 & -0.000241013 & 12 & 12 \\
cer & 50000 & background & $Q_w$ & 24070 & 6.54954e-05 & 0.000202347 & 0.0301229 & 0.0156504 & 0.251474 & 0.251683 & 0.000830166 & 10 & 10 \\
cer & 50000 & background & ThrustMaj & 24070 & 0.000592115 & 9.41716e-05 & 0.036599 & 0.0238961 & 0.371494 & 0.37144 & -0.00014408 & 5 & 5 \\
cer & 100000 & all & $\tau_1^{\mathrm{WTA}}$ & 100000 & 7.69385e-05 & 2.57641e-05 & 0.0197975 & 0.0122643 & 0.239835 & 0.239808 & -0.000113761 & 11 & 11 \\
cer & 100000 & all & $\tau_2^{\mathrm{WTA}}$ & 100000 & 4.49714e-05 & -3.64333e-05 & 0.023961 & 0.0161316 & 0.289812 & 0.289637 & -0.000604453 & 6 & 6 \\
cer & 100000 & all & $\tau_3^{\mathrm{WTA}}$ & 100000 & -2.48832e-05 & 7.33463e-05 & 0.0252972 & 0.0152809 & 0.266812 & 0.266816 & 1.5526e-05 & 9 & 9 \\
cer & 100000 & all & $\tau_4^{\mathrm{WTA}}$ & 100000 & 0.00111636 & 0.00103918 & 0.0586034 & 0.0407436 & 0.913843 & 0.913847 & 3.91345e-06 & 1 & 1 \\
cer & 100000 & all & $\sqrt{d_{12}}$ & 100000 & -0.000610889 & -0.000222564 & 0.0316628 & 0.0192531 & 0.264029 & 0.264112 & 0.000313906 & 10 & 10 \\
cer & 100000 & all & $\sqrt{d_{23}}$ & 100000 & -0.000671922 & -0.000355348 & 0.0377278 & 0.0231882 & 0.381687 & 0.380972 & -0.00187455 & 5 & 5 \\
cer & 100000 & all & ECF$_1$ & 100000 & 0.00099967 & 0.000358012 & 0.0274771 & 0.0136107 & 0.153926 & 0.154481 & 0.00360355 & 12 & 12 \\
cer & 100000 & all & ECF$_2$ & 100000 & 0.000171945 & -5.82356e-06 & 0.0109505 & 0.00616603 & 0.0442861 & 0.0444161 & 0.00293499 & 15 & 15 \\
cer & 100000 & all & ECF$_3$ & 100000 & -0.00437206 & -0.00108583 & 0.0617407 & 0.0371051 & 0.708505 & 0.70588 & -0.00370589 & 2 & 2 \\
cer & 100000 & all & $C_2$ & 100000 & 0.000136103 & 7.97752e-05 & 0.0107965 & 0.00650782 & 0.0960965 & 0.0959975 & -0.00103118 & 14 & 14 \\
cer & 100000 & all & $D_2$ & 100000 & 0.00021342 & 3.8527e-05 & 0.0240104 & 0.0138465 & 0.285772 & 0.2858 & 9.95943e-05 & 7 & 7 \\
cer & 100000 & all & $L_2$ & 100000 & -0.00133374 & -0.000734501 & 0.0366764 & 0.024738 & 0.412908 & 0.412193 & -0.00173 & 3 & 3 \\
cer & 100000 & all & $L_3$ & 100000 & -9.63784e-05 & 1.11759e-08 & 0.0199739 & 0.0117524 & 0.137674 & 0.137922 & 0.00179496 & 13 & 13 \\
cer & 100000 & all & $Q_w$ & 100000 & -0.00115803 & -0.000327706 & 0.0325667 & 0.0177252 & 0.271597 & 0.270842 & -0.00278023 & 8 & 8 \\
cer & 100000 & all & ThrustMaj & 100000 & -0.000983684 & -0.000457016 & 0.0391804 & 0.0263127 & 0.410739 & 0.409584 & -0.00281183 & 4 & 4 \\
cer & 100000 & signal & $\tau_1^{\mathrm{WTA}}$ & 51859 & -0.000110159 & 3.55393e-05 & 0.0185117 & 0.0114388 & 0.199527 & 0.199586 & 0.000299701 & 10 & 10 \\
cer & 100000 & signal & $\tau_2^{\mathrm{WTA}}$ & 51859 & -0.000437295 & -0.00024116 & 0.0233586 & 0.0162019 & 0.259082 & 0.259138 & 0.000215912 & 8 & 8 \\
cer & 100000 & signal & $\tau_3^{\mathrm{WTA}}$ & 51859 & -0.000360852 & -0.000111759 & 0.0270019 & 0.0164175 & 0.296127 & 0.296084 & -0.000147136 & 6 & 6 \\
cer & 100000 & signal & $\tau_4^{\mathrm{WTA}}$ & 51859 & 0.000598283 & 0.000595808 & 0.0573848 & 0.0406569 & 0.87508 & 0.875504 & 0.000485103 & 1 & 1 \\
cer & 100000 & signal & $\sqrt{d_{12}}$ & 51859 & -0.00112158 & -0.000502361 & 0.0279876 & 0.0169855 & 0.197473 & 0.197164 & -0.00156532 & 11 & 11 \\
cer & 100000 & signal & $\sqrt{d_{23}}$ & 51859 & -0.00168402 & -0.000931561 & 0.0357995 & 0.0229609 & 0.351449 & 0.350178 & -0.00361741 & 5 & 5 \\
cer & 100000 & signal & ECF$_1$ & 51859 & 0.00020109 & 0.000201166 & 0.0231424 & 0.0128012 & 0.145607 & 0.146358 & 0.0051559 & 12 & 12 \\
cer & 100000 & signal & ECF$_2$ & 51859 & 0.000173929 & -1.98353e-05 & 0.00886352 & 0.00530188 & 0.0342052 & 0.0344253 & 0.00643495 & 15 & 15 \\
cer & 100000 & signal & ECF$_3$ & 51859 & -0.00785092 & -0.00208998 & 0.0672213 & 0.0384979 & 0.625352 & 0.620138 & -0.00833747 & 2 & 2 \\
cer & 100000 & signal & $C_2$ & 51859 & 5.31464e-08 & 5.84871e-05 & 0.00945272 & 0.0058731 & 0.0794096 & 0.0793443 & -0.000822185 & 14 & 14 \\
cer & 100000 & signal & $D_2$ & 51859 & -0.000176278 & -4.57261e-05 & 0.0194656 & 0.0119233 & 0.207659 & 0.20792 & 0.00125713 & 9 & 9 \\
cer & 100000 & signal & $L_2$ & 51859 & -0.00273015 & -0.00152744 & 0.0380728 & 0.0261984 & 0.407673 & 0.406212 & -0.0035839 & 4 & 4 \\
cer & 100000 & signal & $L_3$ & 51859 & -0.000434441 & -9.18433e-05 & 0.0146859 & 0.00905939 & 0.082159 & 0.0825421 & 0.0046632 & 13 & 13 \\
cer & 100000 & signal & $Q_w$ & 51859 & -0.00219434 & -0.000742391 & 0.0351444 & 0.0201252 & 0.28985 & 0.288348 & -0.00518284 & 7 & 7 \\
cer & 100000 & signal & ThrustMaj & 51859 & -0.00229703 & -0.00102568 & 0.041846 & 0.0288827 & 0.447544 & 0.445534 & -0.00448982 & 3 & 3 \\
cer & 100000 & background & $\tau_1^{\mathrm{WTA}}$ & 48141 & 0.000278485 & 1.42306e-05 & 0.0211827 & 0.0132126 & 0.283259 & 0.283134 & -0.000440419 & 9 & 9 \\
cer & 100000 & background & $\tau_2^{\mathrm{WTA}}$ & 48141 & 0.000564484 & 0.00015983 & 0.0246099 & 0.01606 & 0.322914 & 0.322496 & -0.00129458 & 8 & 8 \\
cer & 100000 & background & $\tau_3^{\mathrm{WTA}}$ & 48141 & 0.000337033 & 0.000216484 & 0.0234609 & 0.0141137 & 0.235234 & 0.23529 & 0.000239322 & 11 & 11 \\
cer & 100000 & background & $\tau_4^{\mathrm{WTA}}$ & 48141 & 0.00167446 & 0.0014921 & 0.0599162 & 0.0408385 & 0.955594 & 0.955135 & -0.00047966 & 1 & 1 \\
cer & 100000 & background & $\sqrt{d_{12}}$ & 48141 & -6.07547e-05 & 0.000100762 & 0.035622 & 0.0221969 & 0.335721 & 0.33623 & 0.00151728 & 7 & 7 \\
cer & 100000 & background & $\sqrt{d_{23}}$ & 48141 & 0.000418345 & 0.000218093 & 0.039805 & 0.0234137 & 0.41426 & 0.414145 & -0.000277405 & 4 & 4 \\
cer & 100000 & background & ECF$_1$ & 48141 & 0.00185993 & 0.000543959 & 0.0321465 & 0.014623 & 0.162889 & 0.163231 & 0.00210396 & 13 & 13 \\
cer & 100000 & background & ECF$_2$ & 48141 & 0.000169808 & 1.28262e-05 & 0.0131986 & 0.00737221 & 0.0551456 & 0.0551782 & 0.000589811 & 15 & 15 \\
cer & 100000 & background & ECF$_3$ & 48141 & -0.000624525 & -0.000148237 & 0.0558368 & 0.0357076 & 0.798075 & 0.798241 & 0.000207626 & 2 & 2 \\
cer & 100000 & background & $C_2$ & 48141 & 0.000282661 & 0.000100821 & 0.0122441 & 0.0073145 & 0.114073 & 0.113935 & -0.0012136 & 14 & 14 \\
cer & 100000 & background & $D_2$ & 48141 & 0.000633215 & 0.000163287 & 0.0289063 & 0.0167862 & 0.369919 & 0.369694 & -0.000609792 & 6 & 6 \\
cer & 100000 & background & $L_2$ & 48141 & 0.000170528 & -2.08318e-05 & 0.035172 & 0.0232022 & 0.418551 & 0.41864 & 0.000213255 & 3 & 3 \\
cer & 100000 & background & $L_3$ & 48141 & 0.000267794 & 0.000159144 & 0.0256704 & 0.0160285 & 0.197478 & 0.197578 & 0.000503149 & 12 & 12 \\
cer & 100000 & background & $Q_w$ & 48141 & -4.16886e-05 & 5.04553e-05 & 0.0297899 & 0.0155084 & 0.251934 & 0.251992 & 0.000227361 & 10 & 10 \\
cer & 100000 & background & ThrustMaj & 48141 & 0.000431095 & 4.40031e-05 & 0.036309 & 0.0238565 & 0.371087 & 0.370861 & -0.000608837 & 5 & 5 \\
\end{longtable}
\end{landscape}


\clearpage
% Generated by tables/make_supplementary_tables.py from the frozen
% artifacts. Every value is copied verbatim; see PROVENANCE.md.
% Requires \usepackage{booktabs}.
\begin{table}[!ht]
\centering
\small
\caption{Score-level and rank-level summaries of the paired SHAP-Delta campaign. Each row is one systematic, sample size and class subset. Score changes are in probability space. $\rho$ is the Spearman correlation between the nominal and varied global feature-importance orderings within that subset; its minimum over all twelve systematic--subset combinations at $100{,}000$ pairs is $0.986$, and over the four all-jet combinations $0.993$.}
\label{tab:supp-shap-scores}
\begin{tabular}{l r l r rrr r}
\toprule
Systematic & Size & Subset & $n$ & Mean $\Delta$score & Median $\Delta$score & Mean $|\Delta$score$|$ & $\rho$ \\
\midrule
esup & 25000 & all & 25000 & 0.00126594 & 0.000565549 & 0.0281417 & 0.996429 \\
esup & 25000 & signal & 12961 & -0.00304732 & -0.00116599 & 0.0321121 & 0.982143 \\
esup & 25000 & background & 12039 & 0.00590954 & 0.00141467 & 0.0238672 & 0.992857 \\
esup & 50000 & all & 50000 & 0.00131986 & 0.000569494 & 0.0281489 & 1 \\
esup & 50000 & signal & 25927 & -0.00288043 & -0.00108558 & 0.0320287 & 0.985714 \\
esup & 50000 & background & 24073 & 0.00584364 & 0.00139111 & 0.0239703 & 0.992857 \\
esup & 100000 & all & 100000 & 0.00116452 & 0.000552509 & 0.0279912 & 1 \\
esup & 100000 & signal & 51851 & -0.00288173 & -0.00102192 & 0.0319719 & 0.985714 \\
esup & 100000 & background & 48149 & 0.00552187 & 0.00131438 & 0.0237045 & 0.996429 \\
esdown & 25000 & all & 25000 & -0.00804751 & -0.0021117 & 0.0286176 & 0.989286 \\
esdown & 25000 & signal & 12963 & -0.0110571 & -0.00437677 & 0.0333686 & 1 \\
esdown & 25000 & background & 12037 & -0.00480642 & -0.00106848 & 0.0235012 & 0.996429 \\
esdown & 50000 & all & 50000 & -0.00805625 & -0.0021484 & 0.0290564 & 0.989286 \\
esdown & 50000 & signal & 25927 & -0.0113878 & -0.00467843 & 0.0338766 & 1 \\
esdown & 50000 & background & 24073 & -0.0044681 & -0.00107664 & 0.0238649 & 0.996429 \\
esdown & 100000 & all & 100000 & -0.00805153 & -0.00214166 & 0.0291714 & 0.992857 \\
esdown & 100000 & signal & 51853 & -0.0113177 & -0.00456274 & 0.0341224 & 1 \\
esdown & 100000 & background & 48147 & -0.004534 & -0.00109002 & 0.0238394 & 0.996429 \\
cpos & 25000 & all & 25000 & 1.68394e-05 & 0 & 0.00420883 & 1 \\
cpos & 25000 & signal & 12963 & -2.23581e-05 & 0 & 0.00416909 & 1 \\
cpos & 25000 & background & 12037 & 5.90523e-05 & 0 & 0.00425163 & 1 \\
cpos & 50000 & all & 50000 & 5.47692e-05 & 0 & 0.00420677 & 1 \\
cpos & 50000 & signal & 25926 & -4.53028e-05 & 0 & 0.00409154 & 1 \\
cpos & 50000 & background & 24074 & 0.00016254 & 0 & 0.00433087 & 1 \\
cpos & 100000 & all & 100000 & 3.68065e-05 & 0 & 0.00416695 & 1 \\
cpos & 100000 & signal & 51852 & -4.37585e-05 & 0 & 0.00407881 & 1 \\
cpos & 100000 & background & 48148 & 0.000123569 & 0 & 0.00426188 & 1 \\
cer & 25000 & all & 25000 & -0.000331825 & -6.64443e-05 & 0.0211876 & 1 \\
cer & 25000 & signal & 12964 & -0.0011782 & -0.000375271 & 0.0224901 & 0.996429 \\
cer & 25000 & background & 12036 & 0.000579811 & 6.04306e-05 & 0.0197848 & 1 \\
cer & 50000 & all & 50000 & -0.000223641 & -3.82895e-05 & 0.021367 & 1 \\
cer & 50000 & signal & 25930 & -0.00128461 & -0.000351995 & 0.0226897 & 1 \\
cer & 50000 & background & 24070 & 0.000919317 & 8.38981e-05 & 0.0199422 & 1 \\
cer & 100000 & all & 100000 & -0.000464717 & -6.22502e-05 & 0.0211959 & 1 \\
cer & 100000 & signal & 51859 & -0.00154458 & -0.000404835 & 0.0224571 & 1 \\
cer & 100000 & background & 48141 & 0.000698547 & 7.55626e-05 & 0.0198373 & 1 \\
\bottomrule
\end{tabular}
\end{table}


\end{document}
